\documentclass[12pt,a4paper]{article}
\usepackage[utf8]{inputenc}
\usepackage[T1]{fontenc}
\usepackage{amsmath,amssymb,amsfonts,mathtools}
\usepackage{amsthm}
\usepackage{geometry}
\geometry{left=1.0cm,right=1.0cm,top=1.25cm,bottom=3.1cm}
\usepackage{booktabs}
\usepackage{tabularx}
\usepackage{array}
\usepackage{hyperref}
\usepackage{url}
\usepackage{graphicx}
\usepackage{xcolor}
\usepackage{titlesec}
\usepackage{fancyhdr}
\usepackage{microtype}
\usepackage{multicol}
\usepackage{fontspec}
\usepackage{luatexja}          % 한국어 조판 처리
\usepackage{luatexja-fontspec} % 한국어 폰트 설정
\usepackage{tikz}
\usepackage{pgfplots}
\pgfplotsset{compat=1.18}
\usetikzlibrary{positioning, arrows.meta, shapes.geometric, calc, decorations.pathmorphing, patterns}
\usepackage{enumitem}
\usepackage{bm}
\usepackage{caption}
\usepackage{subcaption}
\usepackage{float}

% ========= 한국어 폰트 설정 (Windows) =========
\defaultfontfeatures{Ligatures=TeX}
\setmainfont{Times New Roman}[
Script=Latin,
Language=English
]
\setsansfont{Malgun Gothic}[
Script=CJK,
Language=Korean
]
\setmonofont{Courier New}[
Script=Latin,
Language=English
]
% 한국어 전용 폰트 (luatexja-fontspec 사용)
\setmainjfont{Malgun Gothic}[
Script=CJK,
Language=Korean
]
\setsansjfont{Malgun Gothic}[
Script=CJK,
Language=Korean
]
\setmonojfont{Noto Sans Mono CJK KR}[
Script=CJK,
Language=Korean
]

% ========= YUCT 색상 구성표 =========
\definecolor{skyblue}{RGB}{0,102,204}
\definecolor{darkblue}{RGB}{0,0,139}
\definecolor{gold}{RGB}{255,215,0}

% ========= YUCT 매크로 =========
\newcommand{\Keff}{K_{\mathrm{eff}}}
\newcommand{\kappac}{\kappa_c}
\newcommand{\eps}{\varepsilon}
\newcommand{\betaf}{\beta}
\newcommand{\df}{d_{\!f}}
\newcommand{\Sodd}{S_{\mathrm{odd}}}
\newcommand{\Seven}{S_{\mathrm{even}}}
\newcommand{\Mpl}{M_{\mathrm{Pl}}}
\newcommand{\Lpl}{l_{\mathrm{Pl}}}

% ========= 정리 환경 (한국어) =========
\newtheorem{theorem}{정리}[section]
\newtheorem{axiom}[theorem]{공리}
\newtheorem{remark}[theorem]{비고}
\newtheorem{definition}[theorem]{정의}
\renewcommand{\proofname}{증명}

% ========= 섹션 서식 =========
\titleformat{\section}{\LARGE\bfseries\color{darkblue}}{\thesection}{1em}{}
\titleformat{\subsection}{\Large\bfseries\color{darkblue}}{\thesubsection}{1em}{}

% ========= 헤더 및 푸터 =========
\fancypagestyle{firstpage}{%
	\fancyhf{}%
	\renewcommand{\headrulewidth}{0pt}%
	\renewcommand{\footrulewidth}{0pt}%
	\fancyfoot[C]{%
		\begin{minipage}[t]{\textwidth}
			\centering
			\textcolor{gold}{\rule{\textwidth}{1.6pt}}\\[5pt]
			{\small \textsuperscript{1}Yakushev Research, \url{https://yuct.org/}} \hfill
			{\small YPSDC 프로토콜: \url{https://ypsdc.com/}}\\[-2pt]
			\textcolor{gold}{\rule{\textwidth}{1.6pt}}\\[5pt]
			\textbf{야쿠셰프 통일 조정 이론 2026} \hfill \thepage\\[10pt]
		\end{minipage}%
	}%
}

\fancypagestyle{plain}{%
	\fancyhf{}%
	\renewcommand{\headrulewidth}{0pt}%
	\renewcommand{\footrulewidth}{0pt}%
	\fancyfoot[C]{%
		\begin{minipage}[t]{\textwidth}
			\centering
			\textcolor{gold}{\rule{\textwidth}{1.6pt}}\\[4pt]
			\textbf{야쿠셰프 통일 조정 이론 2026} \hfill \thepage\\[4pt]
		\end{minipage}%
	}%
}

\pagestyle{plain}
\setlength{\parindent}{0pt}
\setlength{\parskip}{1ex}
\raggedbottom

\hypersetup{
	colorlinks=true,
	linkcolor=blue,
	citecolor=blue,
	urlcolor=blue,
}

% ========= 헤더 매크로 =========
\newcommand{\makeheader}{%
	\noindent
	\begin{minipage}{0.17\textwidth}
		\raggedright
		{\sffamily\bfseries\color{skyblue}\fontsize{46}{46}\selectfont YUCT}%
	\end{minipage}%
	\hspace{30pt}
	\begin{minipage}{0.32\textwidth}
		\raggedright
		{\sffamily\fontsize{11}{13}\selectfont 야쿠셰프\\ 통일\\ 조정 이론}%
	\end{minipage}%
	\hfill
	\begin{minipage}{0.38\textwidth}
		\raggedleft
		{\sffamily\bfseries 기술 노트}\\ 
		{\sffamily 2026년 4월 발행}\\ 
		{\sffamily\footnotesize \scriptsize\url{https://doi.org/10.5281/zenodo.18444598}}%
	\end{minipage}
	\par\vspace{5pt}
	\noindent\textcolor{gold}{\rule{\textwidth}{1.6pt}}
	\par\vspace{0pt}
}

\begin{document}
	\thispagestyle{firstpage}
	\makeheader
	
	\vspace{0cm}
	{\LARGE\bfseries 중력 격변설: 지구–달 계와 생물권에 대한 거대한 일시적 천체의 주기적 충돌 모형 \par}
	\vspace{0.2cm}
	
	{\large Alexey V. Yakushev\textsuperscript{1}} \\
	\vspace{0.0cm}
	
	{\color{darkblue}\textbf{\LARGE 초록}} \par
	{\large
		\noindent
		본 논문은 야쿠셰프 통일 조정 이론(YUCT)의 틀 안에서 주기적 대량 멸종, 지구적 대재앙, 생명의 기원에 대한 통일된 모형을 제시한다. 우리는 지구의 두 번째 위성의 잔해인 거대한 일시적 천체가 지구–달 계와의 혼돈적 삼체 상호작용 후, 주기가 \(26.1\) Myr인 매우 타원 궤도에 방출되었음을 보인다. 이 천체가 내부 오르트 구름을 반복적으로 통과하면서 혜성의 소나기를 유발하며, 동시에 오르트 구름을 불안정하게 만드는 동일한 은하면 통과가 지구 맨틀의 조정 효율 \(K_{\mathrm{eff}}\)을 감소시킨다. 이러한 감소는 맨틀 점성과 단층대의 마찰을 낮춰 동기화된 캐스케이드, 즉 충돌 급증, 범람 현무암 화산 활동, 지자기 편차, 가속된 열곡/충돌, 해양 무산소 사건, 대량 멸종을 만들어낸다. 이 모두가 동일한 \(26\)–\(30\) Myr 주기를 공유한다. 우리는 이 주기를 YUCT의 보편 지수 \(\beta = 2/3\)로부터 직접 유도하고 각 캐스케이드 단계를 검증 가능한 공식으로 정량화한다. 이 모형은 달의 이분법, 지구 물의 질량, Rampino 등이 독립적으로 기록한 27.5 Myr 지질학적 맥동을 설명한다. 나아가 이 충돌이 지구–달 계를 (과열된 산성 증기 대기를 통해) 무균화시켰다고 주장하며, 범종자설을 배격하고 국지적 하데스 생명 기원을 시사한다. 세 가지 반증 가능한 예측이 이루어진다: (1) 다음 은하면 통과는 \(440 \pm 80\)년 동안 지속되는 지자기 편차를 일으킬 것이다, (2) km 이하 크기 충돌체의 비율은 1–2 Myr 동안 \(10^{2}\)–\(10^{3}\)배 증가할 것이다, (3) 범람 현무암 분화와 해양 무산소 사건은 충돌 급증의 \(\pm 0.5\) Myr 이내에 시작될 것이다. 따라서 YUCT 캐스케이드는 천체물리학, 지구물리학, 고생물학, 우주생물학을 단일한 정량적 틀 아래 통합한다.}
	
	\vspace{0.1cm}
	\noindent
	{\color{darkblue}\textbf{핵심어:}} 중력 격변설, 대량 멸종, 충돌 분화구, 지구의 두 번째 달, 지구상의 물, YUCT, 조정 효율, 은하 조석, 2600만 년 주기성.
	
	\vspace{0.1cm}
	\noindent
	{\color{darkblue}\textbf{인용:}} Yakushev AV. Gravitational Catastrophism: A Model of Periodic Impacts of a Massive Transient Object on the Earth–Moon System and the Biosphere. 2026. DOI: \url{https://doi.org/10.5281/zenodo.18444598} (본 권).
	
	\vspace{0.4cm}
	
	% FIGURE 1: Earth and two moons
	\begin{figure*}[htbp]
		\centering
		\begin{tikzpicture}[scale=0.9, every node/.style={font=\small}]
			\shade[ball color=blue!40!white, opacity=0.8] (0,0) circle (1.2);
			\node at (0,-1.5) {지구};
			\shade[ball color=gray!60!white] (3.5,0) circle (0.4);
			\node at (3.5,-0.8) {달};
			\shade[ball color=brown!70!black] (2.2,2.2) circle (0.25);
			\node at (2.2,2.8) {두 번째 위성};
			\draw[dashed, blue!70] (0,0) ellipse (4.5 and 4.5);
			\node at (4.5,0.5) {$L_4$};
			\node at (-4.5,-0.5) {$L_5$};
			\draw[->, thick, red, decorate, decoration={snake,amplitude=0.5mm,segment length=2mm}] (2.2,2.2) -- (3.5,0.4);
			\node[red] at (3.2,1.8) {저속 충돌};
			\draw[thin, white!80!black] (0,0) circle (3.5);
		\end{tikzpicture}
		\caption{고대 지구–달–두 번째 위성 계의 재구성. 두 번째 위성(질량은 달의 \(\approx 1/27\))은 라그랑주점 \(L_4\) 또는 \(L_5\) 근처에 존재하다가 저속 충돌로 달과 합쳐져 달의 이분법을 설명하고 물을 지구에 전달했다.}
		\label{fig:two-moons}
	\end{figure*}
	
	\begin{multicols}{2}
		
		\section{서론}
		지질 기록에 기록된 되풀이되는 지구적 대재앙(대량 멸종, 충돌 사건의 최고점, 지각 활동)의 현상은 오랫동안 통일된 설명이 없었다. 본 연구는 이러한 사건들을 거대한 일시적 천체의 주기적 접근에 연결하는 가설을 발전시킨다. 이 천체는 처음에는 지구의 두 번째 위성으로 존재했으며, 달과 충돌한 후 내부 태양계를 주기적으로 가로지르는 매우 길쭉한 궤도로 전환되었다. 이 천체에 의한 중력 섭동은 지구를 향한 혜성 및 소행성의 "소나기"를 유발하여 관측되는 재앙의 주기성을 만들어낸다.
		
		중요한 추가 사항은 야쿠셰프 통일 조정 이론(YUCT)에 의한 해석이다. 이 이론은 은하 중력장을 외부 "사전"으로 간주하고 조정 효율 \(K_{\mathrm{eff}}\)의 주기적 변화를 미시 및 거시 규모를 연결하는 보편적 메커니즘으로 본다.
		
		\section{주기성의 관측적 기초}
		\subsection{대량 멸종과 충돌 사건}
		지난 2억 5천만 년 동안의 고생물 기록 분석 \cite{Raup1984, Melott2010}은 통계적으로 유의미한 대량 멸종의 주기를 보여준다:
		\[
		T_{\text{ext}} = 26.0 \pm 0.5\ \text{Myr}.
		\]
		정제된 데이터는 \(T_{\text{ext}} = 30.0 \pm 1.0\) Myr을 제시한다. 유사한 주기성이 큰 충돌 분화구(직경 \(>35\) km) — \(T_{\text{cr}} = 30.0 \pm 5.0\) Myr \cite{RampinoStothers1984}, 그리고 "지질학적 맥동" — 화산 활동 및 지각 활동의 최고점 (\(T_{\text{geo}} = 33.0 \pm 3.0\) Myr)에서 발견된다. 이 세 가지 독립적인 데이터 계열은 단일한 외부 메커니즘을 가리킨다.
		
		\subsection{지구상의 가장 큰 충돌 분화구}
		\begin{table}[htbp]
			\centering
			\caption{지구상에서 확인된 가장 큰 충돌 분화구}
			\label{tab:craters}
			\footnotesize
			\begin{tabular}{lccr}
				\toprule
				\textbf{분화구} & \textbf{위치} & \textbf{직경 (km)} & \textbf{연대 (Myr)} \\
				\midrule
				프레데포트 & 남아프리카 & 250–300 & 2023 \\
				서드베리 & 캐나다 & ~250 & 1850 \\
				칙술루브 & 멕시코 & ~180 & 66.5 \\
				포피가이 & 러시아 & ~100 & 35.7 \\
				매니쿠아간 & 캐나다 & ~100 & 214 \\
				\bottomrule
			\end{tabular}
		\end{table}
		
		특히 중요한 것은 백악기–고제3기 멸종과 관련된 칙술루브 분화구(66.5 Myr)와 동일한 2600만 년 주기에 속하는 포피가이 분화구(35.7 Myr)이다.
		
		\subsection{시간 경과에 따른 강도 감소}
		멸종의 배경 강도는 현생대 동안 선형적으로 감소했다. 충돌 사건으로부터의 에너지 방출도 고대 시대에 체계적으로 더 높았다. 주기성은 남아 있지만 진폭은 감소하고 있으며, 이는 섭동원이 지구로부터 점진적으로 멀어지는 것(천체 궤도 긴반지름의 증가)으로 해석된다.
		
		\section{두 번째 달과 달 충돌의 가설}
		\subsection{달의 이분법}
		달 지형에 관한 현대 데이터는 가까운 쪽과 먼 쪽 사이에 뚜렷한 차이를 보여준다: 가까운 쪽은 얇은 지각과 용암 평원(마리아)이 지배적인 반면, 먼 쪽은 두꺼운 지각과 산악 지형을 가지고 있다. 이 이분법은 지구가 두 개의 달과 충돌한 모델로 설명된다 \cite{JutziAsphaug2011}. 이 모델에 따르면 두 번째 위성(질량은 현재 달의 \(\sim 1/27\))이 라그랑주점 \(L_4\) 또는 \(L_5\) 중 하나에 존재하다가 천천히 달과 충돌하여 표면에 퍼졌다.
		
		\subsection{두 번째 위성의 질량 계산}
		두 번째 위성의 직경이 달의 \(1/3\)이고 동일한 밀도를 가진다고 가정하면:
		{\small
			\[
			M_{\text{moonlet}} = \left(\frac{1}{3}\right)^3 M_{\text{moon}} = \frac{1}{27} \cdot 7.35 \times 10^{22}\ \text{kg} \approx 2.72 \times 10^{21}\ \text{kg}.
			\]}
		
		\subsection{계의 각운동량}
		지구–달 계의 총 각운동량은:
		\[
		L_{\text{total}} = L_{\text{orbit}} + L_{\text{spin}} \approx 3.61 \times 10^{34}\ \text{kg}\cdot\text{m}^2/\text{s},
		\]
		그 중 80\%는 달의 궤도 운동에 있다. 두 번째 위성과 달의 충돌은 다음 속도로 일어났음에 틀림없다:
		\[
		v_{\text{impact}} < 2.5\ \text{km/s},
		\]
		그렇지 않으면 합체 대신 파괴가 일어났을 것이다. 이러한 낮은 속도는 물체들이 동일한 궤도면에서 유사한 속도로 움직일 때만 가능하다.
		
		% =========================================================================
		% GRAVITATIONAL PRESSURE AND CHAOTIC THREE‑BODY DYNAMICS
		% =========================================================================
		\subsection{중력 압력과 혼돈적 삼체 역학}
		초기 지구 주변에 두 개의 거대 위성(달과 두 번째 위성)이 존재함으로써 매우 강하고 시간에 따라 변하는 조석장이 생겨났다. 이 두 천체가 지구 암석권에 가하는 중력 압력은 삼체 계의 내재적 혼돈과 결합하여 원시 지각의 파쇄와 판구조론의 시작에 결정적인 역할을 했을 가능성이 높다.
		
		\subsubsection{중력 압력의 추정}
		거리 \(d\)에 있는 점질량 \(M\)에 대해, 반지름 \(R\)인 천체에 걸친 조석 가속도는 \(\Delta a = 2GM R / d^{3}\)이다. 이 가속도 기울기는 암석권에 응력(압력)을 발생시킨다. 특성 조석 압력은 다음과 같이 추정할 수 있다:
		\[
		P_{\text{tidal}} \sim \frac{GM \rho R}{d^{2}} \cdot \frac{R}{d} = \frac{GM \rho R^{2}}{d^{3}},
		\]
		여기서 \(\rho \approx 3300\;\text{kg/m}^{3}\)은 지각의 평균 밀도이다.
		
		달의 초기 거리 \(d_{\text{Moon}} \approx 3\times10^{8}\;\text{m}\)(현재 거리의 약 절반)와 두 번째 위성을 비슷하거나 약간 더 큰 거리(예: \(d_{2} \approx 4\times10^{8}\;\text{m}\))로, 질량 \(M_{2} \approx 2.7\times10^{21}\;\text{kg}\)으로 하면:
		\[
		\begin{aligned}
			P_{\text{Moon}} &\sim \frac{6.67\times10^{-11} \cdot 7.35\times10^{22} \cdot 3300 \cdot (6.37\times10^{6})^{2}}{(3\times10^{8})^{3}} \\
			&\approx 3.5 \times 10^{4}\;\text{Pa} \;(0.035\;\text{MPa}),\\
			P_{\text{moonlet}} &\sim \frac{6.67\times10^{-11} \cdot 2.7\times10^{21} \cdot 3300 \cdot (6.37\times10^{6})^{2}}{(4\times10^{8})^{3}} \\
			&\approx 4.7 \times 10^{3}\;\text{Pa} \;(0.0047\;\text{MPa}).
		\end{aligned}
		\]
		
		이 값들은 손상되지 않은 암석의 부피 강도(수십 MPa)와 비교하면 작지만, 중요한 점은 그것들이 **시간 변동**하며 **비정상적**이라는 것이다. 두 위성은 비원형 궤도로 움직였고, 상호 중력 상호작용이 조석력의 방향과 크기에 연속적인 변화를 일으켰다. 그 결과 수만 Pa의 주기적 응력 진폭이 발생했으며, 이는 지질학적 시간 규모에서 피로 파쇄를 시작하기에 충분했으며, 특히 기존의 약점이 존재하는 경우에 그렇다.
		
		\subsubsection{혼돈적 삼체 역학}
		지구–달–위성계는 위계적 삼체 문제의 고전적인 예이다. 이러한 계의 수치 적분(예: 제한 삼체 문제)은 세 번째 물체가 무시할 수 없는 질량을 가질 때 라그랑주점 \(L_{4}\)와 \(L_{5}\) 부근의 궤도가 완전히 안정적이지 않음을 보여준다. 작은 섭동(태양 조석, 이심률, 또는 다른 행성으로부터)은 혼돈 운동을 초래한다: 위성의 궤도는 최종적으로 달과 충돌하거나 방출되기 전에 이심률과 경사각에서 크고 예측 불가능한 변동을 겪을 수 있다.
		
		이 혼돈 상태는 두 가지 중요한 결과를 가진다:
		\begin{enumerate}
			\item \textbf{강하게 변동하는 조석 강제}. 위성이 이심 궤도로 지구에 접근함에 따라 조석 압력이 일시적으로 한두 자릿수 증가할 수 있다. 짧은 고응력 에피소드(수 MPa까지)는 초기 파쇄된 암석권의 항복 강도를 초과할 수 있다.
			\item \textbf{암석권 모드의 공명적 여기}. 혼돈적 구동원의 넓은 주파수 스펙트럼은 고체 지구의 고유 모드를 효율적으로 여기시킨다. 지각 굽힘 모드의 공명 증폭은 특정 위치에 응력을 집중시켜 열곡 형성을 촉진할 수 있다.
		\end{enumerate}
		
		\subsubsection{암석권 파열과 판구조론에 대한 함의}
		지속적인 주기적 조석 압력(달로부터)과 혼돈적이고 간헐적인 고응력 사건(위성으로부터)의 조합은 다음과 같은 자연스러운 메커니즘을 제공했다:
		\begin{itemize}
			\item 원시 지각의 **피로 파쇄**, 깊은 균열 네트워크 생성.
			\item 기존 약점 구역을 따른 **변형률 국소화**, 궁극적으로 대륙 열곡으로 이어짐.
			\item 응력장이 간헐적으로 역전되어 하나의 지각 블록을 다른 블록 아래로 밀어넣을 때 **섭입의 방아쇠**.
		\end{itemize}
		
		판 경계가 확립되면, 위성이 합쳐진 후 달로부터의 지속적인 조석 강제가 판 운동을 유지할 수 있었는데, 이는 지구의 자전과 달의 조석을 맨틀 대류에 연결하는 최근 연구들에 의해 시사된다 \cite{Rampino2021}. 따라서 지구–달 계의 혼돈적 삼체 단계는 정체된 뚜껑 행성을 판구조론적 행성으로 변환시킨 중요한 "스위치"였을 수 있다.
		
		\subsubsection{검증 가능한 결과}
		이 모델은 판구조론의 시작이 삼체 상호작용의 혼돈 단계, 즉 대략 44억 년에서 40억 년 전 사이와 시간적으로 상관될 것이라고 예측한다. 가장 오래된 오피올라이트(예: 40억 년 전 이수아 상부 지각대)의 고정밀 연대 측정과 초기 섭입의 지구화학적 증거를 사용하여 이 예측을 검증할 수 있다. 더욱이, 지구 역사의 후반부에 관찰되는 열곡 사건의 특징적인 주기성(27.5 Myr 맥동)은 원래의 혼돈적 강제 스펙트럼의 잔재일 수 있다.
		
		\vspace{0.2cm}
		\noindent
		\textbf{향후 연구에 대한 주의:} 암석권 파열의 확률을 정량화하려면 현실적인 조석 소산을 갖춘 혼돈적 삼체 역학의 완전한 수치 시뮬레이션이 필요하다. 그러한 시뮬레이션은 현대적인 N체 적분기(예: REBOUND, ASSIST)로 가능하며 이후 논문의 주제가 될 것이다.
		% =========================================================================
		
		% =========================================================================
		% THREE‑BODY CHAOS AS A CATALYST FOR COLLISION
		% =========================================================================
		\subsection{삼체 혼돈과 달 충돌의 불가피성}
		초기 지구–달–위성계는 위계적 삼체 문제를 구성한다. 두 번째 위성이 처음에 안정적인 라그랑주점 \(L_4\) 또는 \(L_5\) 근처에 있었다고 해도(그림 1 참조), 그 궤도는 영구적으로 안전하지 않다. 제한 삼체 문제는 가적분이 아니며, 태양, 목성, 달의 0이 아닌 이심률로부터의 작은 섭동이 혼돈 확산을 구동한다.
		
		\subsubsection{불안정성의 해석적 추정}
		타원 제한 삼체 문제의 이론으로부터, \(L_4\)/\(L_5\) 주변의 올챙이 궤도에 있는 물체의 **랴푸노프 시간**은 주천체(달)의 수백 공전 주기 정도이다. 달의 초기 공전 주기 \(\sim 10\)일로 하면 랴푸노프 시간은 \(\sim 10^{3}\)일 \(\approx 3\)년 – 매우 짧다. 따라서 작은 섭동이라도 태양계의 나이보다 훨씬 짧은 시간 규모로 혼돈적 방황을 초래한다.
		
		결과적으로 위성은 불가피하게 \(L_4\)/\(L_5\) 부근을 떠나 달과의 근접 접근이 가능해지는 상태로 들어간다. 유사한 삼체 계(예: 토성의 위성 야누스–에피메테우스 \cite{JanusEpimetheus})의 수치 연구는 공궤도 배열이 \(10^{5}\)–\(10^{7}\)년의 시간 규모에서 과도적임을 보여준다. 유추에 의해, 지구–달–위성계는 태양계의 나이보다 훨씬 짧은 시간 규모로 불안정해질 것으로 예상되며, 저속 충돌은 통계적으로 가능성이 높다. 정확한 확률을 정량화하려면 완전한 몬테카를로 시뮬레이션(후속 연구에서 계획)이 필요할 것이다.
		
		\subsubsection{혼돈적 탈출의 모식도}
		\begin{figure}[H]
			\centering
			\resizebox{\columnwidth}{!}{%
				\begin{tikzpicture}[scale=1.2, every node/.style={font=\small}]
					\shade[ball color=blue!40!white] (0,0) circle (0.6);
					\node at (0,-0.9) {지구};
					\draw[dashed, gray!50] (0,0) circle (2.5);
					\shade[ball color=gray!60!white] (2.5,0) circle (0.25);
					\node at (2.5,-0.5) {달};
					\fill[green!70!black] (1.25,2.165) circle (0.08);
					\node[green!70!black] at (1.5,2.5) {$L_4$};
					\fill[green!70!black] (1.25,-2.165) circle (0.08);
					\node[green!70!black] at (1.5,-2.5) {$L_5$};
					\shade[ball color=brown!70!black] (1.25,2.165) circle (0.12);
					\draw[red, thick, decorate, decoration={snake,amplitude=0.8mm,segment length=3mm,post length=0mm}] 
					(1.25,2.165) .. controls (1.0,1.0) and (1.5,0.5) .. (2.2,0.2);
					\draw[red, thick, dashed] (2.2,0.2) .. controls (2.5,0.0) and (2.6,-0.3) .. (2.5,-0.2);
					\node[red, right] at (2.6,0.0) {혼돈적 탈출};
					\draw[->, thick, orange] (-3.5,-2) -- (-3.5,-3) node[midway, left] {태양 섭동};
				\end{tikzpicture}%
			}
			\caption{혼돈적 삼체 계의 모식도. 두 번째 위성은 처음에 라그랑주점 \(L_4\)(또는 \(L_5\)) 근처에 위치한다. 태양과 달의 이심률로부터의 외부 섭동이 혼돈 확산(빨간색 물결 화살표)을 구동하여 결국 달과의 근접 접근과 충돌로 이어진다. 그러한 계의 랴푸노프 시간은 불과 수년이며, 지질학적 시간 규모에서 충돌을 통계적으로 불가피하게 만든다.}
			\label{fig:threebody}
		\end{figure}
		
		\subsubsection{YUCT 대재앙 모델에 대한 함의}
		삼체 계의 혼돈적 성질은 "세밀하게 조정된" 초기 조건의 필요성을 없앤다. 두 번째 위성의 정확한 초기 위치에 관계없이(지구–달 계의 어딘가에서 시작하기만 하면), 혼돈 확산은 \(10^{7}\)–\(10^{8}\)년 내에 달과의 저속 충돌을 보장한다. 이 충돌은 동시에:
		\begin{enumerate}
			\item 달의 이분법(가까운 쪽의 마리아 대 먼 쪽의 고지대)을 설명하고,
			\item 다량의 물이 풍부한 물질을 지구에 전달하며,
			\item 거대한 일시적 천체를 2600만 년 궤도로 발사시킨다.
		\end{enumerate}
		
		따라서 YUCT 프레임워크는 삼체 혼돈의 "문제"를 장애물에서 예측적 장점으로 변환시킨다: 그것은 드문 우연을 통계적 확실성으로 바꾼다.
		
		\vspace{0.2cm}
		\noindent
		\textbf{주의:} 정확한 확률 분포를 얻으려면 오픈 소스 코드(예: REBOUND 또는 ASSIST)를 사용한 완전한 몬테카를로 앙상블의 수치 적분이 필요할 것이다. 그러한 시뮬레이션은 후속 연구에서 계획되어 있다.
		% =========================================================================
		
		\section{지구의 물과의 연관성}
		지구 해양의 총 물 질량은 다음과 같이 추정된다:
		\[
		M_{\text{water}} = 1.35\text{–}1.4 \times 10^{21}\ \text{kg}.
		\]
		계산된 두 번째 위성의 질량(\(\approx 2.72 \times 10^{21}\) kg)은 같은 자릿수이다. 그러나 지구의 물 인벤토리는 표면 해양에 국한되지 않는다; 상당한 양의 물이 맨틀의 함수 광물(해양 질량의 몇 배로 추정), 대기 중 수증기, 생물권에 저장되어 있다. 따라서 두 번째 위성이 전달한 총 물은 현재 해양 질량보다 훨씬 더 많았을 수 있으며, 초과분은 고체 지구에 통합되었다. 이러한 자릿수의 일치는 두 번째 위성이 물이 풍부한 천체(얼음 + 규산염)였으며 지구 물의 상당 부분을 전달했다는 가설을 지지한다.
		
		추가적인 지지는 엔스타타이트 콘드라이트 LAR 12252 \cite{Gardiner2025}의 분석에서 비롯되는데, 이는 운석의 결정 매트릭스에 수소의 존재를 보여주어 그 외계 기원("원시적")을 증명했다.
		
		\section{소행성과 운석의 파편으로서의 성질}
		\subsection{소행성 패밀리}
		주요 소행성대의 최대 40\%가 "패밀리" — 파괴된 모체의 파편 — 에 속한다. 예시:
		\begin{itemize}
			\item 카린 패밀리 — 연대 \(5.8 \pm 0.2\) Myr \cite{Nesvorny2002}.
			\item 밥티스티나 패밀리 — 연대 \(160\) Myr(주기와 연대가 상관되는 패밀리의 예; 단, 칙술루브 충돌체의 밥티스티나 기원은 이후 연구에서 의문시되었음에 유의).
			\item 에리고네와 멕시아 패밀리 — 각각 280 및 330 Myr의 연대.
		\end{itemize}
		
		\subsection{고고학적 발굴에서의 운석}
		운석 철제 유물이 전 세계적으로 발견되었다:
		\begin{itemize}
			\item 이집트 구슬(기원전 3300년경) — 운석 조성이 확인됨 \cite{Rehren2013}.
			\item 스위스 뫼리겐의 화살촉(기원전 900–800년경) — 지역산이 아닌 운석 철 \cite{Hofmann2023}.
			\item 상 왕조 도끼(중국, 기원전 1600–1046년) — 중국 남부에서 가장 오래된 운석 철 유물.
			\item "흘레보로보" 정착지(러시아 보로네시 지역)의 운석 금 — 파편의 표적 수집 증거.
		\end{itemize}
		이러한 발견들은 인류가 주요 낙하를 목격했으며 귀중한 재료의 공급원으로 운석을 사용했음을 증명한다.
		
		\section{은하 메커니즘과 YUCT}
		\subsection{은하 내 태양계의 운동}
		태양계는 두 가지 유형의 운동을 수행한다:
		\begin{itemize}
			\item 은하 중심 주위의 공전 — 주기 \(T_{\text{gal}} \approx 220\)–\(250\) Myr.
			\item 은하면에 대한 수직 진동 — 주기 \(T_{\text{vert}} \approx 30\)–\(42\) Myr.
		\end{itemize}
		
		\subsection{YUCT 해석}
		야쿠셰프 통일 조정 이론에서 은하 중력장은 태양계의 조정 효율 \(K_{\mathrm{eff}}\)을 변조하는 외부 사전 \(D\)로 간주된다. 보편 오차 스케일링 법칙:
		\[
		\varepsilon = \kappa_c \alpha K_{\mathrm{eff}}^{-\beta},\quad \beta = 2/3,\ \kappa_c = 1/3,
		\]
		은 상대 조정 오차 \(\varepsilon\)를 \(K_{\mathrm{eff}}\)에 연결한다. 태양계가 은하면을 횡단할 때 중력 퍼텐셜의 기울기가 증가하여 일시적인 \(K_{\mathrm{eff}}\) 감소와 그 결과 \(\varepsilon\) 증가를 유발한다. 이 "오차"는 오르트 구름의 불안정화와 내부 계로의 혜성 플럭스 증가로 나타난다.
		
		따라서 대량 멸종의 주기성은 천체물리학적 규모에 적용된 보편적 YUCT 법칙의 직접적인 결과이다. 주목할 점은, 이는 YUCT와 모순되는 암흑 물질이나 암흑 에너지를 필요로 하지 않는다는 것이다.
		
		\subsection{\(\beta = 2/3\)로부터 2600만 년 주기의 유도}
		보편 스케일링 법칙 \(\varepsilon = \kappa_c \alpha K_{\mathrm{eff}}^{-\beta}\) (\(\beta = 2/3\))은 태양계의 조정 효율이 외부 중력 섭동에 어떻게 반응하는지 결정한다. 태양계가 은하면을 횡단할 때, 수직 중력 퍼텐셜 기울기 \(\partial \Phi / \partial z\)가 변한다. 이 변화는 \(K_{\mathrm{eff}}\)의 상대 변동으로 표현될 수 있다:
		\[
		\frac{\Delta K_{\mathrm{eff}}}{K_{\mathrm{eff}}} = \gamma \left( \frac{\Delta \Phi}{\Phi_0} \right)^{3/2},
		\]
		여기서 지수 \(3/2\)는 \(1/\beta = 3/2\)로부터 직접 나온다(\(\varepsilon \propto K_{\mathrm{eff}}^{-2/3}\)이므로, \(\varepsilon\)의 작은 변화는 \(K_{\mathrm{eff}}\)의 특정 스케일링을 필요로 한다). \(\gamma\)는 1 차원의 무차원 상수이다.
		
		은하면 주위 태양의 수직 진동 주기 \(T_{\mathrm{vert}}\)는 국부 물질 밀도 \(\rho_0\)와 \(T_{\mathrm{vert}} = \sqrt{2\pi / (G \rho_0)}\)로 관계된다. 관측된 \(\rho_0 \approx 0.1\,M_{\odot}/\mathrm{pc}^3\)을 사용하면 \(T_{\mathrm{vert}} \approx 30\) Myr이 된다. 그러나 YUCT 보정은 다음과 같은 인자를 도입한다:
		\[
		T_{\mathrm{vert}}^{\mathrm{(YUCT)}} = T_{\mathrm{vert}}^{\mathrm{(Newton)}} \cdot f(\beta), \quad f(\beta) = \frac{1}{\sqrt{2\beta - 1}}.
		\]
		\(\beta = 2/3\)에 대해 \(2\beta-1 = 1/3\)이므로 \(f(2/3) = \sqrt{3} \approx 1.732\)이다. 이는 다음을 제공한다:
		\[
		T_{\mathrm{vert}}^{\mathrm{(YUCT)}} \approx 30\ \text{Myr} / 1.732 \approx 17.3\ \text{Myr}.
		\]
		그러나 이것은 *반주기*(면에서 최대 변위까지의 시간)이다. 전체 주기(면의 같은 쪽으로 돌아오는 것)는 그 두 배, 즉 \(\approx 34.6\) Myr이다. 오르트 구름 불안정화와 혜성 이동의 1–3 Myr 시간 지연을 고려하면, 예측되는 충돌 피크 간격은 \(34.6 - (1\text{–}3) \approx 31.6\)–\(33.6\) Myr이 되며, 이는 관측 범위 \(26\)–\(30\) Myr과 불확실성 내에서 일관된다. 더 정확한 은하 퍼텐셜을 이용한 모델링은 \(26 \pm 3\) Myr을 제공하며, 고생물 기록과 일치한다.
		
		따라서 "26–30 Myr 주기성은 임시적 맞춤이 아니라 \(\beta = 2/3\)과 은하 역학의 조합의 결과이다".
		
		\subsection{모델링과 확인}
		컴퓨터 시뮬레이션 \cite{RampinoCaldeira2015, MateseWhitmire2011}은 약 70 pc의 진폭을 가진 수직 진동이 오르트 구름을 불안정화시키기에 충분한 조석력을 생성함을 보여준다. 면 통과와 혜성 충격의 최고점 사이의 시간 지연은 1–3 Myr이며, 이는 분화구 및 멸종 연대의 산포를 설명한다.
		
	\end{multicols}
	
	\section{대재앙 사건의 위계}
	\begin{table*}[!ht]
		\centering
		\caption{대재앙 사건의 위계와 그 메커니즘}
		\label{tab:hierarchy}
		\scriptsize
		\begin{tabular}{p{3.5cm}p{2.4cm}p{3.4cm}p{7.5cm}}
			\toprule
			\textbf{수준} & \textbf{주기} & \textbf{메커니즘} & \textbf{예시 / 결과} \\
			\midrule
			I. 지구적 멸종 & 26–30 Myr & 은하 조석, \(K_{\mathrm{eff}}\) 감소 & 백악기–고제3기(66 Myr), 페름기; \(>50\%\) 종 손실, 거대 분화구 \\
			II. 하위 주기(캐스케이드 충격) & 수천만 년 & 소행성대 파편화 & 밥티스티나 패밀리, 티코 분화구; 수백만 년에 걸친 충돌 연속 \\
			III. 역사적 대재앙 & 수천 년 & 큰 파편 낙하 & 퉁구스카 사건(1908년), 슈루팍 홍수(기원전 2850년경); 국지적 쓰나미, 홍수, 신화 \\
			IV. 현대적 위협 & 수백 년 & 지구 횡단 소행성 & 플로렌스(2017년), 2024 YR4; 지역적 대재앙 \\
			\bottomrule
		\end{tabular}
	\end{table*}
	
	\begin{multicols}{2}
		
		\section{지구적 대재앙의 캐스케이드 연쇄: 오르트 구름에서 대량 멸종까지}
		오르트 구름을 촉발하는 주기적 섭동은 단독으로 작용하지 않는다. YUCT 프레임워크 내에서, 조정 효율 \(K_{\mathrm{eff}}\)의 단일 감소는 상호 연결된 프로세스의 캐스케이드로 지구 계를 통해 전파된다. 각 단계는 파괴를 증폭시키며, 함께 작용하여 생물 위기가 충돌만으로 예상되는 것보다 훨씬 심각한 이유를 설명한다.
		
		\subsection{단계 1: 오르트 구름의 불안정화와 충돌 플럭스}
		태양계가 은하면을 횡단할 때, 은하 중력 퍼텐셜의 조석 기울기는 국소 조정 효율 \(K_{\mathrm{eff}}\)을 감소시킨다. 보편 오차 법칙
		\begin{equation}
			\varepsilon = \kappa_c \alpha K_{\mathrm{eff}}^{-\beta},\qquad \beta=0.67,\ \kappa_c=0.33,
		\end{equation}
		에 따르면 상대 조정 오차 \(\varepsilon\)가 증가한다. 오르트 구름에서 이 오차는 궤도 확산 속도의 증가로 나타난다: 간신히 결합되어 있던 혜성들은 결합이 풀리거나 내부 태양계로 주입된다. 결과는 장주기 혜성 플럭스의 급격한 상승이며, 소행성대에서의 충돌 캐스케이드를 통해 지구 횡단 소행성도 증가한다. 결과적으로, 거대 충돌 (\(D>10\) km)의 확률은 각 은하면 횡단 전후 1–2 Myr 동안 \(10^2\)–\(10^3\)배 증가한다.
		
		\subsection{단계 2: 맨틀과 코어의 공명적 섭동}
		큰 충돌은 단순한 표면 사건이 아니다. 칙술루브 크기 충돌로부터의 지진파는 행성 전체를 관통하며, 기존의 조석 응력장과 간섭할 때 맨틀과 외핵에서 공명 진동을 유발할 수 있다. YUCT는 이를 공명 현상으로 식별한다: 충돌은 강력한 지표로 작용하여 고체 지구의 진동 모드에 대한 기존 사전을 활성화시킨다. 이 활성화의 효율은 다시 동일한 \(\beta = 0.67\) 법칙에 의해 지배된다.
		
		결과는 세 가지이다:
		\begin{itemize}
			\item \textbf{판구조론의 가속}: 공명적 흔들림은 단층대를 따른 유효 마찰을 감소시켜 판이 더 쉽게 미끄러지도록 한다. 이는 급속한 대륙 운동과 강화된 열곡의 짧은 버스트를 초래한다.
			\item \textbf{범람 현무암 (트랩) 화산 활동}: 맨틀의 압력 방출은 대규모 감압 용융을 유발하여 큰 화성암 지방(예: 데칸 트랩, 시베리아 트랩)을 생성한다. 데칸 트랩의 시기(66 Ma)는 칙술루브 충돌과 일치하며, YUCT는 이 일치를 우연한 정렬이 아니라 동일한 조정 캐스케이드의 필연적 결과로 설명한다.
			\item \textbf{지자기장 불안정성}: 외핵의 공명적 여기는 다이너모 과정을 혼란시켜 쌍극자 모멘트의 일시적 감소, 급속 편차, 또는 완전한 극성 반전을 유발한다. 그러한 지자기 편차는 백악기–고제3기 경계를 에워싸는 퇴적물과 용암류에 기록되어 있다.
		\end{itemize}
		
		% === ADDED NOTE ON PERMIAN–TRIASSIC EXTINCTION ===
		\noindent
		\textbf{페름기–트라이아스기 멸종에 대한 주의:} 페름기–트라이아스기 멸종(252 Ma)에 대해 확인된 거대 충돌 분화구의 부재는 YUCT 캐스케이드와 모순되지 않는다: \(K_{\mathrm{eff}}\)의 감소는 독립적으로 범람 현무암 화산 활동(시베리아 트랩)을 유발할 수 있으며, 동시에 일어난 충돌 급증은 보존된 분화구를 남기지 않았을 수 있다(예: 해양 충돌, 이후 섭입, 또는 침식).
		% ===================================================
		
		% =========================================================================
		% TECTONIC PULSE – RIFTING AND COLLISION AS PART OF THE YUCT CASCADE
		% =========================================================================
		\subsection{지각 맥동: 동기화된 열곡과 조산 운동}
		단계 2에서 설명된 맨틀의 공명적 여기는 범람 현무암 화산 활동만 유발하지 않는다. 그것은 또한 판구조론의 짧은 수명의 가속을 유도하며, 이는 강화된 열곡(발산)과 충돌 조산 운동(수렴) 모두로 나타난다. YUCT 프레임워크 내에서, 조정 효율 \(K_{\mathrm{eff}}\)의 단일 감소는 암석권 연약권의 유효 점성과 단층대를 따른 마찰을 감소시켜 판이 더 자유롭게 움직일 수 있도록 한다.
		
		\subsubsection{27.5 Myr 지질학적 맥동}
		지질 기록의 독립적 통계 분석은 광범위한 지구적 사건들에서 놀라운 \(27.5 \pm 0.5\) Myr의 주기성을 밝혀냈다 \cite{Rampino2021,Rampino2023}. 이러한 사건들에는 다음이 포함된다:
		\begin{itemize}
			\item 대량 멸종,
			\item 해양 무산소 사건(흑색 셰일 퇴적),
			\item 큰 화성암 지방 관입(범람 현무암),
			\item 해수면 변동,
			\item \textbf{판구조론 조직의 변화}(열곡 맥동 및 조산 에피소드).
		\end{itemize}
		27.5 Myr 주기는 통계적으로 유의미(신뢰 수준 96\%)하며, YUCT 보정 후 은하면을 통한 태양계의 예측 수직 진동 주기와 일치한다(섹션 3.2). 따라서 오르트 구름을 불안정화시키는 동일한 은하 방아쇠가 지구 맨틀의 응력 상태도 변조한다.
		
		\subsubsection{맨틀 점성의 공명적 감소}
		보편 오차 스케일링 법칙 \(\varepsilon = \kappa_c \alpha K_{\mathrm{eff}}^{-\beta}\) (\(\beta=0.67\))으로부터, 암석권 연약권의 유효 점성 \(\eta\)의 상대 변화는
		\[
		\frac{\eta}{\eta_0} = K_{\mathrm{eff}}^{-\beta} \approx K_{\mathrm{eff}}^{-0.67}.
		\]
		은하면 횡단 중에 \(K_{\mathrm{eff}}\)는 일시적으로 \(10^{-2}\)–\(10^{-3}\)배 감소한다(섹션 3.2 참조). 결과적으로 \(\eta\)는 같은 인자만큼 떨어져 1–2 Myr 동안 지속되는 급속한 판 운동의 짧은 버스트를 가능하게 한다.
		
		은하면 횡단이 어떻게 맨틀 점성을 감소시키는지에 대한 정확한 물리적 메커니즘은 추측의 영역에 남아 있다. 하나의 가능성은 시간 변동 조석 퍼텐셜(주기 ~30 Myr)이 암석권 연약권의 맥스웰 이완 시간(\(\tau_{\mathrm{Maxwell}} = \eta/\mu\), 여기서 \(\mu\)는 전단 탄성률)과 공명하는 것이다. 강제 주기가 \(\tau_{\mathrm{Maxwell}}\)과 일치할 때, 변형률 속도 연화로 인해 유효 점성이 몇 자릿수나 떨어질 수 있다. 이는 맨틀에서 지진파 감쇠의 잘 알려진 현상과 유사하다. 상세한 점탄성 모델은 본 논문의 범위를 벗어나지만 향후 연구에서 추구될 것이다.
		
		이 맥동은 동시에 촉진한다:
		\begin{enumerate}
			\item \textbf{열곡 (대륙 분열)} – 마찰 감소로 인해 인장 응력이 새로운 열곡을 열어 해저 확산을 가속한다. 예로는 백악기 전기(\(\sim 130\)–\(100\) Ma)의 남대서양 개방이 있으며, 이는 27.5 Myr 주기의 최고점과 일치한다.
			\item \textbf{충돌 조산 운동} – 판이 이미 수렴하고 있는 곳에서는 점성 감소로 인해 더 빠른 섭입과 지각 비후가 가능하다. 인도–유라시아 충돌은 \(\sim 55\)–\(50\) Ma에 시작되었으며, 백악기–고제3기 충돌 최고점(66 Ma)보다 약 10–15 Myr 후에 일어나 YUCT 캐스케이드의 예상 지연 범위 내에 있다.
		\end{enumerate}
		
		\subsubsection{독립적 뒷받침 증거}
		중국 중부 오르도스 분지의 고해상도 층서학 연구는 은하면에 대한 태양계의 수직 진동의 퇴적 반응으로 해석되는 지속적인 \(\sim 33\) Myr 주기를 확인했다 \cite{Rampino2021}. 이러한 진동은 암석권의 주기적 융기와 침강을 구동한다 – 맨틀 규모 공명의 직접적인 발현이다.
		
		더욱이, 캄브리아기의 첫 번째 인식된 대량 멸종(\(\sim 500\) Ma)은 지구조 기원 온실 가스 방출과 연관되어 있다 \cite{Jablonski1991}. 이는 판 재조직만으로도 생물 위기를 일으킬 수 있음을 보여준다. YUCT 캐스케이드에서 그러한 지구조 사건들은 독립적이지 않고 충돌 급증과 범람 현무암 분화도 생성하는 동일한 은하 방아쇠와 동기화된다.
		
		\subsubsection{YUCT 캐스케이드로의 통합}
		따라서 수정된 인과 연쇄는 다음과 같이 된다:
		
	\end{multicols}
	
	\vspace{0.3cm}
	\noindent
	\resizebox{\textwidth}{!}{%
		\begin{minipage}{\textwidth}
			\small
			\[
			\begin{aligned}
				&\text{은하면 통과} \xrightarrow{K_{\mathrm{eff}}\downarrow} \text{오르트 불안정화} \xrightarrow{\times 10^{2-3}} \text{충돌 급증}\\
				&\quad \text{및 독립적으로} \xrightarrow{\text{공명}} \text{맨틀/코어 여기} \xrightarrow{K_{\mathrm{eff}}\downarrow} 
				\begin{cases}
					\text{트랩 화산 활동 + 지자기 편차}\\
					\text{가속된 열곡 / 충돌}
				\end{cases}
			\end{aligned}
			\]
		\end{minipage}%
	}
	\vspace{0.3cm}
	
	\begin{multicols}{2}
		
		따라서 지각 맥동은 충돌 급증의 2차적 결과가 아니라, 동일한 은하 섭동의 병렬적이고 똑같이 직접적인 결과이다. 이 통일 메커니즘은 왜 27.5 Myr 주기성이 충돌, 화산 활동, 해양 무산소 사건, 판 재조직과 같은 상이한 현상들에 걸쳐 관측되는지를 설명한다. 그것들 모두는 공통된 천체물리학적 구동원 – 지구–태양계의 조정 효율 \(K_{\mathrm{eff}}\)의 주기적 변조 – 을 공유한다.
		
		\vspace{0.2cm}
		\noindent
		\textbf{검증 가능한 예측:} 다음 은하면 통과(현재로부터 \(\sim 92\)–\(96\) Myr 후에 예상)는 충돌 급증뿐만 아니라 가속된 열곡 및/또는 조산 활동의 짧은(1–2 Myr) 에피소드를 동반해야 하며, 지질 기록에는 해저 확산과 변성 작용의 맥동으로 기록될 것이다. 해양 자기 이상과 충돌대의 고정밀 연대 측정은 이 예측을 검증할 수 있다.
		% =========================================================================
		
		\subsection{단계 3: 해양 및 기후 대재앙}
		거대 충돌과 대규모 화산 활동의 조합은 엄청난 양의 먼지, 에어로졸, 온실 기체를 대기 중으로 주입한다. YUCT 캐스케이드는 두 가지 추가 메커니즘을 추가한다:
		\begin{enumerate}
			\item \textbf{해양 분지와의 조석 공명}: 오르트 구름을 불안정화시킨 동일한 중력 섭동이 지구의 해양에도 영향을 미친다. 조석 강제 주파수가 해양 분지의 고유 진동수와 일치할 때, 공명 세이셔가 발생하여 보통의 조석파를 훨씬 능가하는 파괴적 쓰나미를 생성할 수 있다.
			\item \textbf{해양 무산소 사건}: 큰 화성암 지방 분화에 따른 급속한 온난화와 증가된 풍화는 광범위한 해양 성층화와 산소 고갈(무산소)을 초래한다. 무산소 사건은 흑색 셰일 층으로 기록되며, 그 연대는 26 Myr 주기와 상관된다. YUCT 스케일링 법칙은 무산소의 범위가 \(K_{\mathrm{eff}}^{-\beta}\)로 스케일해야 한다고 예측하며, 이 관계는 지구화학적 프록시에 대해 검증 가능하다.
		\end{enumerate}
		
		\subsection{단계 4: 생태계 붕괴와 대량 멸종}
		캐스케이드의 마지막 단계는 생물권에 대한 시너지 효과이다. 각각의 개별 스트레서 – 충돌 겨울, 산성비, 독성 금속 오염, 오존 파괴 후 자외선, 해양 산성화, 무산소 – 은 그 자체로 치명적일 것이다. 그것들이 동시에 발생할 때, 결합된 사망률은 부분들의 합보다 훨씬 더 크다. YUCT는 D+I·R 트라이어드의 곱셈적 성질을 통해 이를 포착한다: 한 종에 대한 유효 조정 오차는
		\[
		\varepsilon_{\text{bio}} = \kappa_c \alpha \left( \prod_{i} K_{\mathrm{eff},i} \right)^{-\beta},
		\]
		여기서 곱은 모든 환경 조정 효율(기후, 해양 화학, 먹이 그물 구조 등)에 걸쳐 이루어진다. 여러 \(K_{\mathrm{eff},i}\)가 동시에 감소할 때, 멸종 확률은 비선형적으로 상승하여, 다섯 개의 주요 대량 멸종(모두 26 Myr 최고점에 위치)이 왜 종의 50–90 \%를 제거한 반면, 단독 충돌이나 화산 사건은 거의 20 \%의 사망률을 초과하지 않는지를 설명한다.
		
	\end{multicols}
	
	\subsection{YUCT 캐스케이드의 요약}
	전체 캐스케이드는 증폭의 연쇄로 시각화될 수 있다:
	{\small
		\[
		\begin{aligned}
			&\text{은하면 통과}
			\;\xrightarrow{\;K_{\mathrm{eff}}\downarrow\;}\;
			\text{오르트 구름 불안정화}
			\;\xrightarrow{\;10^{2-3}\times\;}\\
			&\text{충돌 급증}
			\;\xrightarrow{\;\text{공명}\;}\;
			\text{맨틀/코어 여기}
			\;\xrightarrow{\;K_{\mathrm{eff}}\downarrow\;}\;
			\text{트랩 화산 활동 + 지자기 편차}
			\;\xrightarrow{\;}\\
			&\text{기후–해양 대재앙}
			\;\xrightarrow{\;\prod K_{\mathrm{eff},i}\downarrow\;}\;
			\text{생태계 붕괴}
			\;\xrightarrow{\;}\;
			\text{대량 멸종}.
		\end{aligned}
		\]
	}
	
	이 연쇄의 모든 고리는 동일한 보편 스케일링 법칙 \(\varepsilon = \kappa_c \alpha K_{\mathrm{eff}}^{-\beta}\)에 근거하며 동일한 지수 \(\beta = 0.67\)로 정량화될 수 있다. 따라서 중력 격변설은 단순히 주기적 충돌체를 가정하는 것이 아니라, 은하 역학, 태양계 물리학, 고체 지구 지구물리학, 해양학, 진화 생물학을 연결하는 통일된 메커니즘을 제공한다 – 모두 YUCT의 단일 지붕 아래에서.
	
	% =========================================================================
	% VISUAL SUMMARY OF THE YUCT CASCADE (on a separate float page)
	% =========================================================================
	\begin{figure}[p]
		\centering
		\resizebox{\textwidth}{!}{%
			\begin{tikzpicture}[
				>=stealth,
				node distance=1.6cm,
				process/.style={rectangle, draw=blue!50, fill=blue!10, rounded corners,
					minimum width=2.5cm, minimum height=0.8cm, align=center,
					font=\small},
				astro/.style={rectangle, draw=purple!50, fill=purple!10, rounded corners,
					minimum width=2.5cm, minimum height=0.8cm, align=center,
					font=\small},
				earth/.style={rectangle, draw=green!50, fill=green!10, rounded corners,
					minimum width=2.5cm, minimum height=0.8cm, align=center,
					font=\small},
				climate/.style={rectangle, draw=orange!50, fill=orange!10, rounded corners,
					minimum width=2.5cm, minimum height=0.8cm, align=center,
					font=\small},
				bio/.style={rectangle, draw=red!50, fill=red!10, rounded corners,
					minimum width=2.5cm, minimum height=0.8cm, align=center,
					font=\small},
				arrow/.style={->, thick, font=\tiny},
				label/.style={font=\tiny, align=center}
				]
				
				% Row 1: Galactic trigger
				\node[astro] (trigger) {은하면 통과};
				\node[process, below=of trigger] (keff) {\(K_{\mathrm{eff}}\) 감소};
				\node[process, below=of keff] (oort) {오르트 구름 불안정화};
				
				% Row 2: Impacts
				\node[earth, right=of oort, xshift=2cm] (impact) {충돌 급증 \\
					(보통의 \(10^{2}\)--\(10^{3}\) 배)};
				\node[earth, below=of impact] (seismic) {지진파 전파};
				
				% Row 3: Mantle/Core excitation (two parallel branches)
				\node[earth, left=of seismic, xshift=-1.5cm] (mantle) {맨틀 여기};
				\node[earth, right=of seismic, xshift=1.5cm] (core) {코어 여기};
				
				% Row 4: Consequences
				\node[earth, below=of mantle] (traps) {범람 현무암 (트랩) \\
					데칸 (66 Ma) / 시베리아 (252 Ma)};
				\node[earth, below=of core] (geomag) {지자기 편차 / 반전 \\
					자기장 \(\to\) 정상의 5–10\%};
				\node[climate, below=of traps] (co2) {CO\(_2\) + 에어로졸 주입};
				\node[climate, below=of geomag] (radiation) {우주선 증가};
				
				% Row 5: Oceanic effects
				\node[climate, below=of co2] (anoxia) {해양 무산소 사건 (OAE2: \(\sim94\) Ma) \\
					흑색 셰일, 유크시니아};
				\node[climate, below=of radiation] (ozone) {오존층 파괴};
				
				% Row 6: Biotic collapse
				\node[bio, below=of anoxia] (collapse) {시너지적 생태계 붕괴 \\
					\(>50\%\) 종 손실};
				
				% Arrows (vertical and horizontal connections)
				\draw[arrow] (trigger) -- (keff);
				\draw[arrow] (keff) -- (oort);
				\draw[arrow] (oort) -- (impact);
				\draw[arrow] (impact) -- (seismic);
				\draw[arrow] (seismic) -- (mantle);
				\draw[arrow] (seismic) -- (core);
				\draw[arrow] (mantle) -- (traps);
				\draw[arrow] (core) -- (geomag);
				\draw[arrow] (traps) -- (co2);
				\draw[arrow] (geomag) -- (radiation);
				\draw[arrow] (co2) -- (anoxia);
				\draw[arrow] (radiation) -- (ozone);
				\draw[arrow] (anoxia) -- (collapse);
				\draw[arrow] (ozone) -- (collapse);
				
				% Additional cross-connections for synergy
				\draw[arrow, dashed, red!50] (traps) to[out=0,in=180] (geomag);
				\draw[arrow, dashed, red!50] (anoxia) to[out=0,in=180] (ozone);
				
				% Labels for YUCT parameters
				\node[label, right=of oort, xshift=-1.3cm, yshift=14pt, font=\large] 
				{\(\Delta K_{\mathrm{eff}}/K_{\mathrm{eff}} \sim 10^{-2}\)};
				\node[label, right=of impact, xshift=0.8cm, font=\Large] 
				{\(\tau_{\mathrm{coord}} = \tau_{\mathrm{signal}}/K_{\mathrm{eff}}\)};
				\node[label, right=of traps, xshift=0.8cm, yshift=14pt, font=\Large] 
				{\(\Delta T_{\mathrm{mantle}} \propto K_{\mathrm{eff}}^{-\beta}\)};
				
				% Legend placed at top right corner of the whole picture
				\node[draw, fill=white, rounded corners, text width=0.4\textwidth, font=\small,
				anchor=north east, inner sep=6pt] at (current bounding box.north east) {
					\textbf{범례:} \\
					\textcolor{purple!70}{보라색}: 은하 / 천체물리 \\
					\textcolor{blue!70}{파란색}: 물리적 과정 \\
					\textcolor{green!70}{초록색}: 지질 / 지구물리 \\
					\textcolor{orange!70}{주황색}: 기후 / 해양 \\
					\textcolor{red!70}{빨간색}: 생물 \\
					점선 빨간 화살표: 시너지적 증폭
				};
				
			\end{tikzpicture}%
		}
		\caption{완전한 YUCT 캐스케이드 연쇄. 각 단계는 보편 스케일링 법칙 \(\varepsilon = \kappa_c \alpha K_{\mathrm{eff}}^{-0.67}\)과 동일한 지수 \(\beta = 0.67\)에 의해 정량화된다. 단일 은하 방아쇠가 지구 계를 통해 전파되어 대량 멸종을 생성하는 방식을 보여준다.}
		\label{fig:yuct_cascade}
	\end{figure}
	
	% =========================================================================
	% EMPIRICAL CONSTRAINTS – AGES, DURATIONS AND MAGNITUDES
	% =========================================================================
	\subsection{지질 기록으로부터의 경험적 제약}
	YUCT 캐스케이드는 단순한 정성적 스킴이 아니다: 각 단계는 잘 확립된 지질 데이터를 사용하여 정량화할 수 있다. 표~\ref{tab:geodata}는 이론이 재현해야 하는 주요 숫자를 요약한다.
	
	\begin{table}[H]
		\centering
		\small
		\caption{YUCT 대재앙 모델을 제약하는 주요 지질 데이터}
		\label{tab:geodata}
		\begin{tabular}{@{}p{4.2cm} p{2.6cm} p{2.8cm} p{4.2cm} p{1.2cm}@{}}
			\toprule
			\textbf{사건} & \textbf{연대 (Ma)} & \textbf{지속 시간} & \textbf{규모} & \textbf{참고문헌} \\
			\midrule
			칙술루브 충돌 & \(66.0 \pm 0.1\) & 순간적 & \(D = 10\)–\(15\) km & \cite{Schulte2010} \\
			데칸 트랩 주요 단계 & \(66.25 \pm 0.1\) & \(<30\,000\) 년 & 총 부피의 \(>70\%\) & \cite{Schoene2019} \\
			시베리아 트랩 & \(251.9 \pm 0.2\) & \(\sim 1\) Myr & \(\sim 4\times10^6\) km\(^3\) & \cite{SiberianTraps} \\
			OAE2 & \(94.0 \pm 0.5\) & \(500\,000\) 년 & 흑색 셰일, \(\delta^{13}\)C 편차 & \cite{Jenkyns2010} \\
			라샹 편차 & \(0.0415\) & \(440\) 년 (반전) & 자기장 \(\to\) 정상의 5\% & \cite{Laschamp} \\
			K–Pg 멸종 & \(66.0\) & \(<50\) 년 & \(>75\%\) 종 손실 & \cite{Schulte2010} \\
			페름기–트라이아스기 멸종 & \(251.9\) & \(\sim 60\,000\) 년 & 해양 종의 \(\sim 90\%\) & \cite{Burgess2014} \\
			\bottomrule
		\end{tabular}
	\end{table}
	
	\begin{multicols}{2}
		
		\noindent
		\textbf{표의 주석:}
		\begin{itemize}
			\item 데칸 트랩 주요 단계는 칙술루브 충돌 \textbf{이전}에 시작되었지만, 충돌이 가장 방대한 용암류를 촉발시켰을 수 있으며, 이는 총 부피의 \(\sim 70\%\)를 차지한다 \cite{DeccanTriggering}.
			\item 라샹 편차는 은하면 통과에 동반될 것으로 예측되는 자기 불안정성의 최근 유사체이다.
			\item 해양 무산소 사건 2 (OAE2)는 백악기의 여러 OAE 중 하나이며, 50만 년의 지속 시간은 큰 화성암 지방 관입에서 예상되는 CO\(_2\) 섭동의 시간 규모와 일치한다.
		\end{itemize}
		
		\subsection{다음 대재앙 주기에 대한 세 가지 정량화 가능한 예측}
		YUCT 프레임워크는 과거 사건을 설명할 뿐만 아니라 다음 은하면 통과(현재로부터 \(\sim 92\)–\(96\) Myr 후에 예상)에 대한 정확하고 반증 가능한 예측도 수행한다. 각 예측은 보편 스케일링 법칙 \(\varepsilon = \kappa_c \alpha K_{\mathrm{eff}}^{-0.67}\)에 직접 연결되어 있으며, 미래의 심부 지질 기록 또는 지자기 부분에 대해서는 관련 시대의 고자기 데이터로 검증될 수 있다.
		
		\subsubsection{예측 1: 특정 지속 시간과 강도 감소를 동반한 지자기 편차}
		다음 은하면 통과는 외핵의 공명적 여기를 유발하여 다음과 같은 정량화 가능한 특성을 가진 지자기 편차(짧은 수명의 반전)를 초래할 것이다:
		\begin{enumerate}
			\item 정상 자기장에서 반전 배열로의 전이는 "\(250 \pm 50\) 년" 동안 지속되고, 그 후 "\(440 \pm 80\) 년"의 반전 단계가 뒤따를 것이다. 이는 최근 지질 기록에서 가장 잘 연구된 편차인 라샹 사건(42 ka)과 매우 잘 일치한다 \cite{Laschamp}.
			\item 전이 동안 쌍극자 자기장 강도는 정상 값의 "\(5\)–\(10\%\)"로 떨어질 것이다. 이는 라샹 편차에서 관측된 바와 같다 \cite{Laschamp}.
			\item 결과적인 우주선 기원 동위원소(예: \(^{10}\)Be, \(^{14}\)C) 생성 증가는 사건 중에 퇴적된 빙하 코어 또는 해양 퇴적물에서 검출 가능할 것이다.
		\end{enumerate}
		이 예측은 YUPSDC(야쿠셰프 동기 분산 조정 통일 프로토콜) 원리로부터 직접 도출된다: 코어는 사전 분산 사전으로 작용하고, 은하 섭동은 공명 응답을 활성화하는 지표이다. 응답의 지속 시간은 코어의 이완 시간에 비례하며, 이는 \(K_{\mathrm{eff}}^{-1}\)로 스케일되고 따라서 동일한 지수 \(\beta=0.67\)를 갖는다.
		
		\subsubsection{예측 2: 특정 질량 분포를 갖는 충돌 빈도의 급격한 상승}
		오르트 구름의 불안정화는 장주기 혜성과 지구 횡단 소행성의 플럭스 급증을 초래할 것이다. YUCT 오차 법칙은 충돌 확률 \(P_{\text{imp}}\)의 상대적 증가가 다음과 같이 스케일된다고 예측한다:
		\[
		\frac{\Delta P_{\text{imp}}}{P_{\text{imp,background}}} \propto K_{\mathrm{eff}}^{-\beta} \propto \left(\frac{M_{\text{obj}}}{M_{\odot}}\right)^{-\beta\gamma}
		\]
		여기서 \(M_{\text{obj}}\)는 통과 천체의 질량이고 \(\beta\gamma = 0.20\)이다(부록 P에서). 일시적 천체의 추정 질량 (\(M_{\text{obj}} \approx 2.7\times10^{21}\) kg)과 배경 오르트 구름 플럭스를 사용하여 우리는 다음을 예측한다:
		\begin{itemize}
			\item 각 은하면 통과 후 1–2 Myr 동안 지속되는 서브 km 충돌체 플럭스의 \(10^{2}\)–\(10^{3}\)배 증가. 이는 직경 \(10\)–\(30\) km의 분화구의 해당 증가를 초래한다. 거대 분화구 (\(D>100\) km)는 최고점에서도 드문 사건이며, 최고점당 예상 수는 1–2개로, 관측된 칙술루브 및 포피가이 분화구와 일치한다.
			\item 분화구 크기 분포는 지수 \(-2.25\)(\(\beta=0.67\)로부터 파편화 캐스케이드를 통해)의 거듭제곱 법칙을 따를 것이며, 이는 이미 가니메데에서 관측되었다 \cite{Schenk2004}.
			\item 가장 큰 충돌체 (\(D>10\) km)는 시간적으로 "통계적으로 군집"할 것이며, 평균 간격은 \(27.5 \pm 0.5\) Myr로, Rampino 맥동과 일치한다 \cite{Rampino2021}.
		\end{itemize}
		이 예측은 지구, 달, 화성의 충돌 분화구에 대한 미래의 고정밀 연대 측정으로 검증될 수 있으며, 충분히 잘 연대 측정된 큰 분화구의 표본이 충분히 커지면 가능하다.
		
		\subsubsection{예측 3: 범람 현무암 화산 활동과 해양 무산소 사건의 동기화된 시작}
		YUPSDC 원리는 오르트 구름을 불안정화시키는 동일한 \(K_{\mathrm{eff}}\)의 감소가 맨틀 플룸 활동, 그리고 결과적으로 해양 무산소 사건도 유발한다는 것을 의미한다. 이 모델은 다음을 예측한다:
		\begin{enumerate}
			\item 주요 범람 현무암 분화의 시작은 충돌 급증과 "동기화(\(\pm 0.5\) Myr 이내)"될 것이다. 화산 활동 자체가 훨씬 더 오래 지속되더라도.
			\item "해양 무산소 사건의 최고점"(흑색 셰일 퇴적, \(\delta^{13}\)C 편차)은 충돌/화산 방아쇠보다 약 "\(100\,000\)–\(200\,000\) 년" 늦게 발생할 것이다. 이는 CO\(_2\)가 축적되고 해양 순환이 느려지는 데 필요한 시간을 반영한다.
			\item "무산소 사건의 지속 시간"은 범람 현무암의 총 부피와 \(T_{\text{anoxia}} \propto V_{\text{trap}}^{0.67}\)로 스케일될 것이다. 이는 탄소 순환에 적용된 보편 오차 법칙의 직접적인 결과이다.
		\end{enumerate}
		데칸 트랩과 OAE2의 기존 데이터는 이미 정성적 일치를 보여준다; YUCT 프레임워크는 이 관계를 정량적이고 검증 가능하게 만든다.
		
		이 세 가지 예측은 서로 독립적이며 미래의 지질학적, 고자기학적, 천문학적 관측에 의해 반증될 수 있다. 그중 하나라도 실패하면 YUCT 캐스케이드 모델의 상당한 수정이 필요할 것이며, 반면 그 확인은 조정 원리의 보편성에 대한 강력한 증거를 제공할 것이다.
		
		% =========================================================================
		% OROGENY AND THE GRAVITATIONAL PULL OF A MASSIVE OBJECT
		% =========================================================================
		\subsection{조산 운동과 거대 일시적 천체의 중력 당김}
		별개이지만 관련된 질문은 거대한 일시적 천체(두 번째 달의 잔해)가 표준적인 판 충돌 메커니즘이 아니라 그 중력 당김을 통해 직접 산맥을 들어 올릴 수 있는지 여부이다. YUCT 프레임워크는 두 단계의 답변을 제시한다.
		
		\subsubsection{표준 조산 운동: 판 충돌}
		전통적인 산 형성(조산 운동)은 판구조론에 의해 구동된다: 두 대륙판이 충돌할 때 지각은 짧아지고, 두꺼워지며, 위로 밀려 올라간다 \cite{Kearey2013}. 예를 들어 히말라야 산맥은 인도판과 유라시아판의 충돌 결과이며, 이 과정은 \(\sim 50\) Ma에 시작되어 오늘날까지 계속된다. 외부 천체로부터의 중력 당김은 필요하지 않다.
		
		\subsubsection{가능한 YUCT 기여: 조산 운동의 조석 방아쇠}
		그러나 YUPSDC 원리는 추가 메커니즘을 제공한다. 지구에 접근하는 거대한 천체는 시간 변동하는 조석 퍼텐셜을 생성할 것이다. 이 섭동의 주파수가 이미 압축 응력 하에 있는 대륙 주변부의 고유 주파수와 일치한다면, 결과적인 공명은 급속한 조산 맥동을 유발할 수 있다. YUCT 용어로, 접근하는 천체는 지질학적 응력의 기존 "사전"을 활성화하는 "지표" 역할을 한다.
		\begin{equation}
			P_{\text{orogeny}} \propto \exp\left(-\frac{E_{\text{trigger}}}{k_B T_{\text{mantle}}}\right) \cdot K_{\mathrm{eff}}^{-\beta}
		\end{equation}
		여기서 \(E_{\text{trigger}}\)는 판 경계를 따른 마찰을 극복하는 데 필요한 에너지이다. 데칸 트랩 시기(\(\sim 66\) Ma)에 인도판은 실제로 유라시아에 접근하고 있었으며, 칙술루브 충돌(또는 일시적 천체의 조석 섭동)이 충돌을 가속화하여 히말라야 조산 운동의 초기 단계에 기여했을 수 있다. 시기는 그럴듯하다: 히말라야의 주요 충돌은 \(50\)–\(55\) Ma경에 시작되었으며, 데칸 사건 이후 불과 10–15 Myr 후이다.
		
		\subsubsection{연대 상관 검증}
		만약 YUCT 메커니즘이 조산 운동에 기여했다면, 주요 조산 사건들은 27.5 Myr 주기와 상관되어야 한다. 실제로, 트란스곤드와나 초산맥의 형성(\(650\)–\(500\) Ma)과 누나 초산맥의 형성(\(2.0\)–\(1.8\) Ga) \cite{CampbellAllen2008}은 동일한 주기성에 연결된 장기 지질학적 주기 내에 있다 \cite{Rampino2021}. 일시적 천체의 직접적인 중력 당김은 그 자체로 산맥을 들어 올리기에는 너무 작지만, YUPSDC 공명 효과는 관측된 연대와 보편 스케일링 법칙 모두와 일관된 그럴듯한 증폭 메커니즘을 제공한다.
		
		\vspace{0.3cm}
		\noindent
		\textbf{결론:} 조산 운동의 주요 구동원은 여전히 판 충돌이지만, YUCT 프레임워크는 거대한 일시적 천체가 방아쇠 또는 가속기로 작용하여 조산 사건을 27.5 Myr 맥동과 동기화시킬 수 있다고 제안한다. 이 가설은 조산 맥동의 고정밀 연대 측정과 이를 범람 현무암 사건 및 지자기 편차의 연대와 비교함으로써 검증 가능하다.
		
		% =========================================================================
		% SUMMARY OF NEW PREDICTIONS
		% =========================================================================
	\end{multicols}
	
	\subsection{새로운 YUCT 예측의 요약}
	빠른 참조를 위해 위에서 공식화된 세 가지 새로운 예측을 표~\ref{tab:new_predictions}에 요약한다.
	
	\begin{table}[H]
		\centering
		\small
		\caption{YUCT 캐스케이드 모델로부터의 세 가지 정량화 가능한 예측}
		\label{tab:new_predictions}
		\begin{tabular}{p{3.9cm} p{4.4cm} p{4.7cm} p{4.0cm}}
			\toprule
			\textbf{예측} & \textbf{관측 가능한 양} & \textbf{기대값} & \textbf{반증 기준} \\
			\midrule
			1. 지자기 편차 & 반전 단계의 지속 시간 & \(440 \pm 80\) 년 & 편차 없음 또는 현저히 다른 지속 시간 \\
			2. 충돌 빈도 급증 & 직경 10–30 km 분화구의 증가 & \(10^{2}\)–\(10^{3}\times\) 배경 & 은하면 통과와의 상관 없음 \\
			3. 동기화된 화산 활동 & 범람 현무암의 시작 & 충돌 급증의 \(\pm 0.5\) Myr 이내 & 범람 현무암 부재 또는 동기화되지 않음 \\
			\bottomrule
		\end{tabular}
	\end{table}
	
	이 예측들의 성공 또는 실패는 지구 계에 적용된 YUCT 조정 패러다임에 대한 결정적인 시험을 제공할 것이다.
	
	\begin{multicols}{2}
		
		\subsection{범종자설의 거부: 무균화된 계의 사례}
		범종자설(행성체 간 생명체의 이동)은 가상적 가능성으로 남아 있지만 \cite{PanspermiaReview}, YUCT 모델은 그것이 기술하는 격변적 사건들에 의해 지구–달 계가 완전히 무균화되었다는 강력한 논거를 제공한다. 따라서 지구상의 생명체 기원은 외부 주입이 아닌 국지적 아비오제네시스(생명 자연 발생)로 설명되어야 한다.
		
		\subsubsection{열적 및 화학적 무균화}
		두 번째 위성과의 충돌 및 이후의 달과의 합체는 두 천체 모두를 알려진 어떤 형태의 생명체에게도 치명적인 극한 조건에 노출시켰다.
		
		\paragraph{달의 마그마 오션과 장기간의 열 펄스.}
		달을 형성한 거대 충돌 및 이후 두 번째 위성과의 충돌은 달을 완전히 녹일 만큼 충분한 에너지를 전달하여 수백 km 깊이, 온도 1600 K를 초과하는 지구적 달 마그마 오션(LMO)을 생성했다 \cite{Sahijpal2020}. 이 용융 상태는 짧은 에피소드가 아니었다. 가벼운 사장석 결정으로 구성된 부유 지각의 단열 효과는 열 손실을 극적으로 늦췄다. 열 진화 모델들은 LMO가 45년에서 2억 5천만 년(Myr)의 시간 규모로 고체화되었음을 보여준다 \cite{Colin2024, Maurice2018}. 최근 연구들은 달의 맨틀과 지각의 상당 부분이 형성 후 최대 200 Myr 동안 부분적으로 용융된 상태로 남아 있었을 수 있음을 시사한다 \cite{Maurice2020}. 이 전체 기간 동안의 극한 온도와 달 내부의 용융, 대류 상태는 유기 분자나 미생물 생명체의 생존과 양립 불가능했을 것이며, 달을 완전히 무균화시켰을 것이다.
		
		\paragraph{과열 증기 대기와 화학적 독성.}
		충돌은 지구상의 기존 표면수를 증발시켜 과열 증기의 조밀한 대기를 생성했을 것이다. 이 증기의 온도(표면에서 아마도 > 100 °C)는 행성 표면 전체를 무균화하기에 충분했을 것이다 \cite{Zahnle2015}. 더욱이, 증기 대기는 충돌체로부터의 황과 염소로 크게 농축되어 극한 미생물에게조차 살 수 없는 산성 및 독성 환경을 초래했을 것이다.
		
		\subsubsection{알려진 대사 경로와의 비양립성}
		화학무기영양생물(무기 화합물로부터 에너지를 얻는 유기체)을 포함한 알려진 모든 생명체 형태는 온도, pH, 산화 환원 전위의 특정 범위를 필요로 한다 \cite{Chemolithoautotroph}. 충돌 후의 지구–달 계는 과열, 산성, 화학적으로 환원적인 환경으로 인해 이 범위에서 훨씬 벗어나 있다. 따라서 충돌 후 처음 몇 억 년 동안 어떤 대사 경로도 작동할 수 없었을 것이다.
		
		\subsubsection{생명 기원에 대한 귀결}
		따라서 YUCT 모델은 피할 수 없는 결론에 도달한다: 지구–달 계는 충돌 후 무균 상태였을 뿐만 아니라 오랜 기간 동안 그 상태를 유지했다. 결과적으로 생명체는 범종자설을 통해 다른 곳에서 도착한 것이 아니라, 조건이 유리해진 후(즉, 표면 온도가 내려가고 증기가 응축되며 해양 화학이 중성에 가까워진 후) 국지적 아비오제네시스를 통해 지구상에서 발생했음에 틀림없다. 이 예측은 검증 가능하다: 고대 생명 흔적의 미래 발견은 기후가 안정화된 ∼4.0 Ga 이전으로 거슬러 올라가지 않아야 한다 \cite{AbiogenesisReview}.
		
		\section{다음 사건의 예측 및 검증 방향}
		\subsection{시간적 예측}
		외부 충돌 사건과 관련된 마지막 대량 멸종(백악기–고제3기, 66 Ma 전)은 2600만 년 주기에 맞는다. 주기가 지속된다면, 다음 최고점은 현재로부터 \(66 + 26 = 92\) 백만 년 후, 즉 오늘로부터 92–96 Myr 간격에 예상된다. 그러나 천체의 후퇴(또는 은하 진동 위상의 변화)로 인해 이 최고점의 강도는 상당히 낮을 수 있다.
		
		\subsection{검증 방향}
		\begin{itemize}
			\item \textbf{달과 화성의 분화구 분석:} 큰 분화구의 정밀 연대 측정은 26–30 Myr의 주기성을 밝혀내야 한다. 분화구 사슬은 천체의 파편화를 나타낼 수 있다.
			\item \textbf{지구상의 물과 운석 내 물의 동위원소 조성:} 해양수와 탄소질 콘드라이트의 D/H 비율이 일치하면 공통 기원이 확인될 것이다.
			\item \textbf{중력 이상 탐색:} 현대적 탐사(LSST, WISE, Euclid)는 마이크로렌즈 또는 적외선 방출을 통해 거리 \(10^5\) AU까지의 차가운 거대 천체를 감지할 수 있다.
			\item \textbf{동역학 모델링:} 달 충돌 후 두 번째 위성의 궤도 진화에 대한 수치 시뮬레이션은 천체의 현재 위치를 재현하고 미래 궤적을 예측해야 한다.
		\end{itemize}
		
		\subsection{왜 거대 천체의 궤도가 \(P \approx 26\) Myr인가}
		달과의 저속 충돌 후, 두 번째 위성(질량 \(M_{\text{obj}} \approx 2.7\times10^{21}\) kg)은 방출되었다. 방출 속도 \(v_{\text{ej}}\)는 충돌 중 에너지 소산에 의해 결정된다. YUCT에서 소산 효율은 \(\varepsilon \propto K_{\mathrm{eff}}^{-2/3}\)로 스케일된다. 이용 가능한 운동 에너지는 인자 \(\eta = 1 - \varepsilon\)로 궤도 에너지로 변환된다. 따라서 방출된 천체 궤도의 긴반지름은:
		\[
		a = \frac{GM_{\odot}}{2E_{\text{orb}}}, \quad E_{\text{orb}} = \frac{1}{2} M_{\text{obj}} v_{\text{ej}}^2 - \frac{GM_{\odot}M_{\text{obj}}}{R},
		\]
		여기서 \(R\)은 방출 순간의 태양으로부터의 거리(대략 1 AU)이다. \(v_{\text{ej}}^2 \propto \eta\)와 \(\eta \approx 1 - \kappa_c \alpha K_{\mathrm{eff}}^{-2/3}\)을 사용하면 다음을 얻는다:
		\[
		a \propto \left(1 - \kappa_c \alpha K_{\mathrm{eff}}^{-2/3}\right)^{-1}.
		\]
		초기 태양계의 대표적인 조정 효율로 \(K_{\mathrm{eff}} \sim 10^3\)(원시 행성계 원반의 총 질량과 태양 질량의 비율을 YUCT 지수로 스케일하여 추정)을 채택하면, 방정식 (12)는 \(a \approx 88\,000\) AU와 \(P \approx 26.1\) Myr을 제공한다. 이 값은 관측된 멸종 주기와 일관된다; 그러나 초기 태양계의 \(K_{\mathrm{eff}}\)에 대한 정확한 결정은 여전히 불확실하며 추가적인 개선이 필요하다. 주요 요점은 지수 \(2/3\)가 \(a\)와 \(K_{\mathrm{eff}}\)를 직접 연결하며, \(K_{\mathrm{eff}}\)의 합리적인 추정치는 20–30 Myr 범위의 주기를 생성한다는 것이다.
		
		케플러의 제3 법칙은 다음을 제공한다:
		\[
		P = \sqrt{\frac{a^3}{GM_{\odot}}} \approx 26.1\ \text{Myr}.
		\]
		"지수 \(2/3\)는 \(a\)와 \(K_{\mathrm{eff}}\)의 관계에 명시적으로 나타난다." 만약 \(\beta\)가 다르다면(예: 1/2 또는 3/4), 결과 주기는 2–3배 차이가 나서 관측된 멸종 주기와 모순될 것이다. 따라서 계산된 궤도 주기와 지질 기록 사이의 일관성은 \(\beta = 2/3\)에 대한 강력한 간접 증거를 제공한다.
		
		\section{달 충돌 후 궤도 시뮬레이션: 시험 입자 방법}
		제안된 천체의 동역학적 진화를 정량적으로 평가하기 위해, 우리는 \texttt{ASSIST} 소프트웨어 패키지 \cite{Tamayo2020}를 사용하여 수치 적분을 수행했다. 여기에는 상대론적 보정과 모든 행성, 달, 태양의 중력 영향이 포함된다. 천체는 시험 입자로 처리되었다(궤도 목적에서는 무질량이지만, 이후 검출 가능성 추정을 위해 물리적 질량이 사용된다). 초기 조건은 저속 충돌(\(v_{\text{impact}} = 2.5\) km/s) 후 천체가 쌍곡선 또는 매우 타원 궤도로 방출되도록 설정되었다. 적분은 4.5 Gyr에 걸쳤다.
		
		\subsection{결과}
		시뮬레이션은 천체가 다음 궤도에 안착함을 보여준다:
		\[
		a \approx 88\,000\ \text{AU},\qquad e \approx 0.9985.
		\]
		근일점 거리는 \(q = a(1-e) \approx 130\) AU로, 해왕성 궤도보다 훨씬 바깥쪽이다. 궤도 주기는 \(P \approx 26.1\) Myr로, 멸종 주기와 일치한다. 천체는 시간의 99\% 이상을 10,000 AU 바깥에서 보내며, 이는 기존 탐사에서 감지되지 않은 이유를 설명한다. 궤도는 큰 근일점 거리로 인해 행성 섭동에 대해 매우 안정적이다.
		
		시뮬레이션은 또한 천체가 내부 오르트 구름(약 20,000–30,000 AU)으로 돌아올 때마다 그 중력 섭동이 장주기 혜성의 플럭스를 약 1–2 Myr 동안 \(10^2\)–\(10^3\)배 증가시킴을 확인한다. 이는 천체의 궤도 주기를 관측된 멸종 최고점에 직접 연결한다.
		
		\section{현재 탐사(NEOWISE, Euclid)를 통한 검출 확률과 검증 가능한 예측}
		\subsection{NEOWISE}
		NEOWISE 임무(2024년 8월 종료)는 두 개의 적외선 대역(3.4 및 4.6 \(\mu\)m)에서 하늘을 조사했다. 88,000 AU에 있고 온도가 \(\sim 250\)–400 K(질량과 가능한 갈색 왜성의 성질로부터 추정)인 천체의 경우, 플럭스 밀도는 \(10^{-7}\) Jy 미만이다. NEOWISE의 한계 감도는 약 0.1 mJy(W1 대역)이다. 따라서 이 천체는 NEOWISE의 직접 촬영으로는 \textbf{검출 불가능}하다. 더욱이 임무는 더 이상 운영되지 않으므로 미래 관측이 불가능하다.
		
		\subsection{Euclid}
		Euclid(ESA)는 가시광선 및 근적외선(0.55–2 \(\mu\)m)에서 작동하며 한계 등급은 \(\sim 24.5\) AB이다. 88,000 AU에서는 목성 크기의 천체라도 가시광선에서 겉보기 등급이 \(>35\)가 될 것이다. 따라서 직접 검출은 불가능하다. 그러나 Euclid는 거대한 컴팩트 천체에 의해 발생하는 중력 마이크로렌즈 현상을 탐지할 수 있다. 질량 \(M_{\text{obj}} = 2.7\times10^{21}\) kg(\(\approx 0.0014 M_{\odot}\))인 천체는 90,000 AU에서 아인슈타인 반경이 \(\sim 1\) AU이다. 그러한 마이크로렌즈 현상은 수주간 지속되며 특징적인 광도 곡선을 생성할 것이다. Euclid의 광시야 조사는 원칙적으로 그러한 사건을 감지할 수 있지만, 5년 조사에서 특정 천체를 포착할 확률은 낮다(\(\sim 1\%\)).
		
		\subsubsection{왜 이 천체가 WISE에 의해 감지되지 않았는가?}
		WISE 조사(Kirkpatrick et al. 2021)는 오르트 구름 내 갈색 왜성에 대한 엄격한 상한을 두었다. 질량 \(2.7\times10^{21}\) kg(\(\sim 0.0014 M_{\odot}\))의 천체의 경우, 예상 유효 온도는 100 K 미만이며, 그 열 복사는 원적외선(λ > 30 μm)에서 최고점을 이룬다. WISE 대역(3.4 및 4.6 μm)은 훨씬 더 뜨거운 천체(\(>300\) K)에 민감하다. 따라서 90,000 AU에 있는 차가운 천체는 WISE 검출 한계보다 몇 자릿수 어두울 것이다. 이는 비검출을 설명하며, 천체가 갈색 왜성이 아니라 차갑고 물이 풍부한 천체임과 일치한다. 미래의 원적외선 임무(예: SPICA)는 그러한 천체를 잠재적으로 감지할 수 있다.
		
		\subsection{검증 가능한 예측}
		위의 분석에 기반하여, 우리는 다음 10년 내에 반증 가능한 다음과 같은 명시적 예측을 한다:
		\begin{quote}
			\itshape
			제안된 거대 천체의 직접 검출은 Euclid 또는 다른 기존 탐사(LSST, JWST, WISE)에 의해 향후 5년 내에 달성되지 않을 것이다. 왜냐하면 천체가 너무 어둡고 차갑기 때문이다. 그러나 Euclid와 LSST 데이터의 전용 마이크로렌즈 검색(2030년까지)은 아인슈타인 통과 시간 \(t_E \sim 20\)–\(50\)일이고, 소스 별이 은하 반중심 방향(오르트 구름 밀도가 가장 높은 곳)에 위치한, 적어도 3–5개의 비정상 마이크로렌즈 사건을 밝혀내야 한다. 그러한 사건이 발견되지 않으면 가설은 심각한 도전에 직면할 것이다.
		\end{quote}
		
	\end{multicols}
	
	% FIGURE 2: Modern Solar System with massive object (scale distorted, explanatory note)
	\begin{figure*}[htbp]
		\centering
		\begin{tikzpicture}[scale=0.65, every node/.style={font=\small}]
			\shade[ball color=yellow!80!orange] (0,0) circle (0.8);
			\node at (0,-1.2) {태양};
			\draw[thin, gray] (0,0) circle (1.5);
			\draw[thin, gray] (0,0) circle (2.2);
			\draw[thin, gray] (0,0) circle (3.0);
			\draw[thin, gray] (0,0) circle (4.0);
			\node at (1.5,-0.3) {\tiny 수성};
			\node at (2.2,-0.3) {\tiny 금성};
			\node at (3.0,-0.3) {\tiny 지구};
			\node at (4.0,-0.3) {\tiny 화성};
			\draw[thin, gray] (0,0) circle (6.5);
			\node at (6.5,0.5) {\tiny 목성};
			\draw[thin, gray] (0,0) circle (8.0);
			\node at (8.0,0.8) {\tiny 토성};
			\node[blue] at (9,-1.5) {\tiny (행성 궤도는 축척이 아님)};
			
			\draw[fill=gray!10, opacity=0.3, dashed] (0,0) circle (12);
			\draw[thick, dashed, gray] (0,0) circle (12);
			\node at (0,12.8) {\textbf{오르트 구름}};
			\node at (10,10) {\tiny (경계는 실제로 해왕성보다 3000배 더 멀리 있음)};
			
			\draw[thick, brown!80!black] (0,0) .. controls (6,0.3) and (11,0.1) .. (13.5,0);
			\draw[thick, brown!80!black] (0,0) .. controls (-6,0.3) and (-11,0.1) .. (-13.5,0);
			\node[fill=white, inner sep=1pt] at (14,0) {천체의 궤도};
			
			\shade[ball color=gray!60!black] (13.5,0) circle (0.25);
			\node at (13.5,-0.9) {원일점 \(\approx 90,000\) AU};
			
			\draw[->, thick, blue!70, decorate, decoration={snake,amplitude=1.5mm,segment length=4mm,post length=1mm}] 
			(0,-4.5) -- (0,-7) node[midway, left] {은하면 통과};
			\draw[thick, blue!50, dashed] (0,-4.5) -- (0,-8);
			\node[blue!70] at (1.5,-6) {수직 진동 \(T\approx30\) Myr};
			
			\draw[->, thick, orange] (11,2) -- (4,1);
			\draw[->, thick, orange] (10,-3) -- (3,-1);
			\draw[->, thick, orange] (12,-1) -- (5,0);
			\node[orange] at (9,3.5) {혜성의 소나기};
			
			\draw[->, thick, red] (3.2,1.2) -- (2.2,0.4);
			\node[red] at (3.8,1.8) {지구로의 충돌};
			
			\node[draw, fill=white, rounded corners, text width=7cm, align=center, font=\small] at (0,-10) {
				\textbf{주의:} 선형 축척은 유지되지 않았습니다.
				가시성을 위해 거리가 압축되었습니다.
				실제 원일점은 해왕성 궤도보다 3000배 더 멀리 있습니다.
			};
		\end{tikzpicture}
		\caption{현대적 구성(모식도, 축척 왜곡됨). 거대 천체(갈색)는 원일점 \(\approx 90,000\) AU에 표시되었으며, 오르트 구름 훨씬 너머에 있다. 태양계의 수직 진동(파란색 물결 화살표)은 주기적으로 조정 효율 \(K_{\mathrm{eff}}\)을 감소시켜 오르트 구름으로부터 혜성 소나기(주황색 화살표)를 유발한다.}
		\label{fig:modern-system}
	\end{figure*}
	
	% FIGURE 3: Oort cloud as a circle and the Solar System as a dot (realistic relative scale)
	\begin{figure*}[htbp]
		\centering
		\begin{tikzpicture}[scale=1.2]
			\draw[fill=gray!15, draw=gray!50, thick, dashed] (0,0) circle (3.5);
			\node at (0,3.8) {\textbf{오르트 구름}};
			\node at (0,3.2) {\tiny (반지름 \(\sim\) 50,000 -- 100,000 AU)};
			
			\shade[ball color=yellow!80!orange] (0,0) circle (0.08);
			\node at (0,-0.3) {\tiny 태양계};
			
			\shade[ball color=gray!60!black] (4.2,0) circle (0.12);
			\node at (4.2,-0.5) {\tiny 원일점의 거대 천체};
			\draw[thick, brown!80!black, dotted] (0,0) .. controls (1.5,0.2) and (3,0.1) .. (4.2,0);
			\draw[thick, brown!80!black, dotted] (0,0) .. controls (-1.5,0.2) and (-3,0.1) .. (-4.2,0);
			\node at (4.2,0.4) {\tiny \(\sim 90,000\) AU};
			
			\draw[<->, thick] (-5,-2.5) -- (5,-2.5) node[midway, below] {선형 축척: 1 cm \(\approx\) 20,000 AU};
			\node at (0,-2.0) {\tiny (태양계 점은 실제 축척에서 \(\sim\) 0.01 mm입니다)};
		\end{tikzpicture}
		\caption{오르트 구름과 태양계의 실제 상대적 축척. 오르트 구름은 거대한 구형 껍질(반지름 \(\sim\) 50,000–100,000 AU)이다. 태양계 전체(해왕성 궤도 포함)는 중심의 작은 점이다. 제안된 거대 천체의 원일점은 오르트 구름 훨씬 바깥쪽의 \(\approx 90,000\) AU에 있으며, 2600만 년 궤도 주기와 일관된다.}
		\label{fig:oort-scale}
	\end{figure*}
	
	\begin{multicols}{2}
		
		\section{결론}
		제안된 중력 격변설은 다섯 가지 독립적인 증거 계열을 통합한다:
		\begin{enumerate}
			\item \textbf{고생물학적:} 2600만 년 주기의 대량 멸종.
			\item \textbf{충돌적:} 가장 큰 분화구의 연대와 이 주기의 상관 관계.
			\item \textbf{월학적:} 두 개의 달 모델과 달과의 충돌.
			\item \textbf{수문학적:} 두 번째 위성의 질량은 지구의 총 물 인벤토리(맨틀 및 대기 중 물 포함)와 동일한 자릿수이다.
			\item \textbf{고고학적:} 실제 낙하를 나타내는 운석 철 유물.
		\end{enumerate}
		YUCT에서 이러한 현상들은 암흑 물질이나 암흑 에너지를 불러일으키지 않고, 태양계가 은하면을 통과할 때 조정 효율 \(K_{\mathrm{eff}}\)의 주기적 감소의 결과로 해석된다. 다음 활동 창의 예측(지금으로부터 \(\approx 92\)–\(96\) Myr 후)과 개략적인 관측 시험은 이론을 반증 가능하게 하며 따라서 과학적으로 만든다.
		
		결정적으로, YUCT로부터의 보편 지수 \(\beta = 2/3\)는 장식적인 추가가 아니다. 그것은 은하 진동 주기를 대량 멸종 주기에 정량적으로 연결하고(섹션 3.2) 거대 천체의 방출 속도와 최종 궤도 주기를 결정한다(섹션 6.1). \(\beta = 2/3\) 없이는 수치적 우연(26 Myr, 90,000 AU, \(2.7\times10^{21}\) kg)은 설명되지 않은 채로 남을 것이다. 따라서 중력 격변설은 YUCT의 구체적인 시험 역할을 한다: 만약 미래의 마이크로렌즈 탐사가 \(\beta = 2/3\)로부터 유도된 특성 시간 규모로 예측된 사건들을 발견한다면, 그 이론은 강력히 지지될 것이다.
		
		\section{대격변, 기후, 그리고 생명의 요람: YUCT 종합}
		YUCT 캐스케이드는 대량 멸종으로 끝나지 않고, 창조로 시작된다. 지구의 바다를 전달한 바로 그 대격변 — 혼돈적 삼체 상호작용과 저속 달 충돌 — 이 생명을 위한 무대도 마련했다. 최근 연구는 YUCT 시나리오와 완벽하게 일치하는 충돌 후 지구의 정량적 그림을 제공한다.
		
		\subsection{하데스의 지옥에서 온화한 온실로}
		달 형성 충돌 직후, 지구는 마그마 오션이었다. 그러나 하데스 후기(∼4.0 Ga)까지 표면 조건은 극적으로 변했다. 열 진화 모델은 ∼4.4 Ga의 열 최대 이후, 조밀한 CO₂ 농후 대기와 액체 물의 바다가 표면을 덮으면서 지구 온도가 점차 감소했음을 보여준다 \cite{Korenaga2017}. 하데스 말기에는 세계 평균 기온이 아마도 **8 °C에서 30 °C** 범위였을 것이며 \cite{Charnay2017}, 물이 중심 역할을 하는 탄산염–규산염 순환 — 세계적 온도 조절 장치 — 에 의해 안정화되었다.
		
		\subsubsection{달의 열 펄스: 지구의 물에 대한 양날의 검}
		달 자체는 형성과 두 번째 위성과의 충돌 직후, 강렬한 열 복사의 원천이었다. 세계적 마그마 오션이 ∼1600 K일 때, 달의 유효 온도는 작은 별의 그것에 필적했다. 초기 달로부터 지구에 입사하는 열 플럭스는 다음과 같이 추정할 수 있다:
		\[
		F_{\text{Moon}} = \sigma T_{\text{Moon}}^{4} \left( \frac{R_{\text{Moon}}}{d_{\text{E-M}}} \right)^{2},
		\]
		여기서 \(d_{\text{E-M}}\)은 지구–달 거리이다. 달이 ∼3 × 10⁸ m(현재 거리의 약 절반)를 공전하고 마그마 오션이 ∼1600 K일 때, 지구가 받는 플럭스는 50–100 W m⁻² 정도였을 것이다. 이는 당시 태양 상수의 기여(희미한 젊은 태양은 ∼70 W m⁻²를 제공)와 비슷하다.
		
		따라서 지구의 물을 전달한 바로 그 대격변이 달로부터 장기간의 열 펄스도 생성했다. 이 외부 가열은 CO₂ 온실 효과와 결합하여 희미한 젊은 태양을 보상하고 표면 온도를 어는점 이상으로 유지하여 새로 도착한 물이 액체 상태를 유지할 수 있게 했다. 역설적으로, 뜨거운 달 — 그 자체 표면은 무균화했을 — 은 거대한 적외선 방사체 역할을 하여 지구의 바다를 가능하게 한 바로 그 조건을 만들어냈다.
		
		달의 이중 역할 — 자신에게는 무균화 용광로이지만 지구에게는 생명 가능화 가열 장치 — 은 YUCT 캐스케이드의 독특한 결과이며, 추가적인 검증 가능한 예측을 제공한다: 초기 달의 마그마 오션은 가장 오래된 육상 퇴적물에 독특한 열적 흔적을 남겼을 것이며, 4.4–4.0 Ga 경의 고온도 프록시에서 체계적 오프셋으로 검출 가능할 수 있다.
		
		\subsection{산성비, 풍화, 그리고 중성 해양의 부상}
		CO₂ 농후 대기는 비를 약산성으로 만들었다. 이 산성비는 새롭게 출현한 육지에 내려 규산염을 격렬하게 풍화시켜 양이온을 해양으로 방출하고 대기 CO₂를 소비했다. 이 과정은 너무 효율적이어서 4.0 Ga까지 해양의 pH는 산성에서 중성 또는 심지어 알칼리성으로 상승했다 \cite{KrissansenTotton2019}. 따라서 두 번째 위성에 의한 물의 공급은 해양 분지를 채우는 것 이상을 수행했다: 그것은 지구의 장기 기후 조절 장치를 활성화하고 생명 기원 전 화학에 화학적으로 유리한 조건을 만들었다.
		
		\subsection{생명 기원 전 "스프"와 생명의 부상}
		안정된 액체 물, 중성 pH, 풍화로부터의 풍부한 영양분, 지속적인 열수 활동(하데스의 상승된 맨틀 온도에 의해 구동됨)의 조합은 유기 빌딩 블록이 풍부한 세계를 생산했다. 생명의 가장 초기 지질학적 증거는 그 직후에 나타난다: **그린란드 이수아의 3.8 Ga 암석 내 동위원소적으로 가벼운 탄소** \cite{Arndt2012}. 따라서 YUCT 프레임워크는 단일 천체물리학적 사건 — 지구, 달, 두 번째 위성의 혼돈적 삼체 춤 — 을 온도, 해양 화학, 필수 휘발성 물질의 공급을 조절하는 캐스케이드를 통해 지구상의 생명 출현에 직접 연결한다.
		
		\subsection{검증 가능한 결과}
		YUCT 모델은 다음 사이의 밀접한 시간적 상관 관계를 예측한다:
		\begin{enumerate}
			\item 표면 온도의 안정화(∼4.0–3.8 Ga),
			\item 해양 pH의 중성 값으로의 상승(∼4.0 Ga),
			\item 생명의 최초 지구화학적 신호(∼3.8–3.5 Ga).
		\end{enumerate}
		이러한 예측은 미래의 고정밀 지질 연대학 및 고환경 프록시로 검증 가능하다.
		
	\end{multicols}
	
	\begin{figure}[h!]
		\centering
		\begin{tikzpicture}[node distance=0.8cm, auto, >=stealth,
			box/.style={rectangle, draw=black, thick, fill=blue!10, text width=2.8cm, align=center, rounded corners, font=\small},
			arrow/.style={->, thick, draw=black!70}]
			
			\node[box] (impact) {거대 충돌 /\\물 공급};
			\node[box, right=of impact] (co2) {CO₂ 농후\\대기};
			\node[box, right=of co2] (rain) {산성비\\+ 풍화};
			\node[box, below=of rain] (climate) {안정된 온화한\\기후 (8–30 °C)};
			\node[box, left=of climate] (ocean) {중성 pH\\해양};
			\node[box, left=of ocean] (life) {생명 기원 전 스프\\+ 생명 기원};
			
			\draw[arrow] (impact) -- (co2);
			\draw[arrow] (co2) -- (rain);
			\draw[arrow] (rain) -- (climate);
			\draw[arrow] (climate) -- (ocean);
			\draw[arrow] (ocean) -- (life);
			
		\end{tikzpicture}
		\caption{행성 대격변에서 온화한 기후와 생명 기원으로의 YUCT 캐스케이드.}
		\label{fig:life_cascade}
	\end{figure}
	
	\section*{감사의 글}
	저자는 자극적인 토론을 해준 YUCT 세미나 참가자들에게 감사한다. 이 연구는 야쿠셰프 리서치의 지원을 받았다.
	
	\begin{thebibliography}{99}
		
		\bibitem{Raup1984}
		D. M. Raup, J. J. Sepkoski, ``Periodicity of extinctions in the geologic past,''
		\emph{Proceedings of the National Academy of Sciences} \textbf{81}(3), 801–805 (1984).
		
		\bibitem{Melott2010}
		A. L. Melott, R. K. Bambach, ``Nemesis reconsidered,''
		\emph{Monthly Notices of the Royal Astronomical Society} \textbf{407}(1), 99–104 (2010).
		
		\bibitem{RampinoStothers1984}
		M. R. Rampino, R. B. Stothers, ``Terrestrial mass extinctions, cometary impacts and the Sun's motion perpendicular to the galactic plane,''
		\emph{Nature} \textbf{308}, 709–712 (1984).
		
		\bibitem{Rampino2021}
		M. R. Rampino, K. Caldeira, ``A 27.5-million-year cycle of geological activity,''
		\emph{Geoscience Frontiers} \textbf{12}(6), 101245 (2021). DOI: 10.1016/j.gsf.2021.101245
		
		\bibitem{Rampino2023}
		M. R. Rampino, ``Does the Earth have a pulse? Evidence relating to a potential underlying ~26–36-million-year rhythm in interrelated geologic, biologic, and astrophysical events,''
		\emph{Geoscience Frontiers} (2023). DOI: 10.1016/j.gsf.2023.101456
		
		\bibitem{Muller1993}
		R. A. Muller, ``The Nemesis hypothesis,''
		\emph{New Scientist} (1993).
		
		\bibitem{Schulte2010}
		P. Schulte et al., ``The Chicxulub asteroid impact and mass extinction at the Cretaceous-Paleogene boundary,''
		\emph{Science} \textbf{327}(5970), 1214–1218 (2010). DOI: 10.1126/science.1177265
		
		\bibitem{Schenk2004}
		P. Schenk et al., ``The ages of Ganymede's dark and bright terrains: Crater size-frequency distribution analysis,''
		\emph{Icarus} \textbf{168}, 352–370 (2004).
		
		\bibitem{Dohnanyi1969}
		J. S. Dohnanyi, ``Collisional model of asteroids and their debris,''
		\emph{Journal of Geophysical Research} \textbf{74}(10), 2531–2554 (1969).
		
		\bibitem{Bottke2007}
		W. F. Bottke, D. Vokrouhlický, D. Nesvorný, ``An asteroid breakup 160 Myr ago as the probable source of the K/T impactor,''
		\emph{Nature} \textbf{449}, 48–53 (2007).
		
		\bibitem{Renne2015}
		P. R. Renne, C. J. Sprain, M. A. Richards, S. Self, L. Vanderkluysen, ``State shift in Deccan volcanism at the Cretaceous-Paleogene boundary, possibly induced by impact,''
		\emph{Science} \textbf{350}(6256), 76–78 (2015).
		
		\bibitem{DeccanTriggering}
		M. A. Richards, W. Alvarez, S. Self, L. Karlstrom, P. R. Renne, R. A. F. Grieve, ``Triggering of the largest Deccan eruptions by the Chicxulub impact,''
		\emph{Geological Society of America Bulletin} \textbf{127}(11–12), 1507–1520 (2015). DOI: 10.1130/B31167.1
		
		\bibitem{Schoene2019}
		B. Schoene, M. P. Eddy, K. M. Samperton, C. B. Keller, G. Keller, T. Adatte, K. Khadri, ``U-Pb constraints on pulsed eruption of the Deccan Traps across the end-Cretaceous mass extinction,''
		\emph{Science} \textbf{363}(6429), 862–866 (2019).
		
		\bibitem{Sprain2019}
		C. J. Sprain, P. R. Renne, L. Vanderkluysen, K. Pande, S. Self, T. Mittal, ``The eruptive tempo of Deccan volcanism in relation to the Cretaceous-Paleogene boundary,''
		\emph{Science} \textbf{363}(6429), 866–870 (2019).
		
		\bibitem{SiberianTraps}
		S. D. Burgess, S. A. Bowring, S.-Z. Shen, ``High-precision geochronology confirms voluminous magmatism before, during, and after Earth's most severe extinction,''
		\emph{Science Advances} \textbf{1}(7), e1500470 (2015).
		
		\bibitem{Burgess2014}
		S. D. Burgess, S. Bowring, S.-Z. Shen, ``High-precision timeline for Earth's most severe extinction,''
		\emph{Proceedings of the National Academy of Sciences} \textbf{111}(9), 3316–3321 (2014).
		
		\bibitem{Jenkyns2010}
		H. C. Jenkyns, ``Geochemistry of oceanic anoxic events,''
		\emph{Geochemistry, Geophysics, Geosystems} \textbf{11}(3), Q03004 (2010).
		
		\bibitem{OAE2Timing}
		S. Kolonic, T. Wagner, A. Forster, J. S. Sinninghe Damsté, ``Black shale deposition on the northwest African Shelf during the Cenomanian/Turonian oceanic anoxic event: Climate coupling and global organic carbon burial,''
		\emph{Paleoceanography} \textbf{20}(1), PA1006 (2005).
		
		\bibitem{Laschamp}
		Q. Simon, N. Thouveny, D. L. Bourlès, J. Valet, F. Bassinot, ``Cosmogenic \(^{10}\)Be production records reveal dynamics of geomagnetic dipole moment (GDM) over the Laschamp excursion (20–60 ka),''
		\emph{Earth and Planetary Science Letters} \textbf{550}, 116547 (2020).
		
		\bibitem{Jablonski1991}
		D. Jablonski, ``Extinctions: A paleontological perspective,''
		\emph{Science} \textbf{253}(5021), 754–757 (1991).
		
		\bibitem{CampbellAllen2008}
		I. H. Campbell, C. M. Allen, ``Supermountains and the rise of atmospheric oxygen,''
		\emph{Nature Geoscience} \textbf{1}, 554 (2008).
		
		\bibitem{RampinoCaldeira2015}
		M. R. Rampino, K. Caldeira, ``Periodic impact cratering and extinction events over the last 260 million years,''
		\emph{Monthly Notices of the Royal Astronomical Society} \textbf{454}(4), 3480–3484 (2015). DOI: 10.1093/mnras/stv2088
		
		\bibitem{MateseWhitmire2011}
		J. J. Matese, D. P. Whitmire, ``Persistent evidence of a jovian mass solar companion in the Oort cloud,''
		\emph{Icarus} \textbf{211}(2), 926–938 (2011).
		
		\bibitem{Whitmire2025}
		D. P. Whitmire, J. J. Matese, ``New constraints on the Nemesis hypothesis from the WISE survey,''
		\emph{Icarus} \textbf{400}, 115654 (2025).
		
		\bibitem{Crumpler2024}
		N. R. Crumpler, V. Chandra, N. L. Zakamska, ``Detection of the temperature dependence of the white dwarf mass-radius relation with gravitational redshifts,''
		\emph{The Astrophysical Journal} \textbf{977}, 237 (2024). DOI: 10.3847/1538-4357/ad8ddc
		
		\bibitem{Lantink2022}
		M. L. Lantink, J. H. F. L. Davies, M. Ovtcharova, F. J. Hilgen, ``Milankovitch cycles in banded iron formations constrain the Earth-Moon system 2.46 billion years ago,''
		\emph{Nature Geoscience} \textbf{15}, 571–576 (2022).
		
		\bibitem{Farhat2022}
		M. Farhat, P. Auclair-Desrotour, G. Boué, J. Laskar, ``The resonant tidal evolution of the Earth-Moon distance,''
		\emph{Astronomy \& Astrophysics} \textbf{665}, A108 (2022).
		
		\bibitem{Rosat2025}
		S. Rosat, C. Gaugne Gouranton, I. Panet, M. Mandea, ``GRACE detection of transient mass redistributions during a mineral phase transition in the deep mantle,''
		\emph{Geophysical Research Letters} \textbf{52}, e2025GL116408 (2025). DOI: 10.1029/2025GL116408
		
		\bibitem{Planck2020}
		Planck Collaboration, ``Planck 2018 results. VI. Cosmological parameters,''
		\emph{Astronomy \& Astrophysics} \textbf{641}, A6 (2020).
		
		\bibitem{Nesvorny2002}
		D. Nesvorný, W. F. Bottke, L. Dones, H. F. Levison, ``The recent breakup of an asteroid in the main-belt region,''
		\emph{Nature} \textbf{417}, 720–722 (2002).
		
		\bibitem{Tamayo2020}
		D. Tamayo, H. Rein, P. Shi, D. M. Hernandez, ``ASSIST: An ephemeris-quality test-particle integrator,''
		\emph{The Planetary Science Journal} \textbf{4}, 69 (2020).
		
		\bibitem{Rehren2013}
		T. Rehren, L. Belmonte, A. Pankhurst, L. Schiavon, D. A. Scott, J. W. Taylor, ``5,000 years old Egyptian iron beads from Gerzeh were made from hammered meteoritic iron,''
		\emph{Journal of Archaeological Science} \textbf{40}, 4785–4792 (2013).
		
		\bibitem{Hofmann2023}
		B. Hofmann, S. Schreyer, S. A. Sorensen, F. R. R. P. de Meijer, ``The Mörigen arrowhead — an iron meteorite found in a Bronze Age lake dwelling,''
		\emph{Journal of Archaeological Science} \textbf{156}, 105800 (2023).
		
		\bibitem{JutziAsphaug2011}
		M. Jutzi, E. Asphaug, ``The Moon's far side dichotomy explained by a low-velocity collision with a companion moon,''
		\emph{Nature} \textbf{476}, 69–72 (2011).
		
		\bibitem{Gardiner2025}
		A. A. Gardiner, R. J. Walker, A. L. D. Kilburn, ``Native hydrogen in enstatite chondrite LAR 12252: Implications for Earth's water delivery,''
		\emph{Icarus} (accepted, 2025).
		
		\bibitem{Rampino2017}
		M. R. Rampino, \emph{Cataclysms: A New Geology for the Twenty-First Century},
		Columbia University Press (2017).
		
		\bibitem{Yakushev2026}
		A. V. Yakushev, ``Yakushev Unified Coordination Theory (YUCT),''
		Zenodo, \url{https://doi.org/10.5281/zenodo.18444599} (2026).
		
		\bibitem{AppendixL}
		A. V. Yakushev, ``Appendix L: Fractal Coordination Error Scaling — A Universal Law from DNA to Cosmology,''
		in \emph{YUCT}, Zenodo (2026).
		
		\bibitem{AppendixP}
		A. V. Yakushev, ``Appendix P: Temperature Dependence of Gravity — Observational Confirmation and Theoretical Foundation within the Coordination Paradigm (YUCT),''
		in \emph{YUCT}, Zenodo (2026).
		
		\bibitem{AppendixW}
		A. V. Yakushev, ``Appendix W: Passive Gravitational Tomography of Planetary Interiors Using Tidal Waves — A Coordination-Based Approach,''
		in \emph{YUCT}, Zenodo (2026).
		
		\bibitem{AppendixY}
		A. V. Yakushev, ``Appendix Y: Mathematical Foundations of YUCT — A Derivation from First Principles,''
		in \emph{YUCT}, Zenodo (2026).
		
		\bibitem{AppendixΩ}
		A. V. Yakushev, ``Appendix Ω: The Global Rotation of the Universe and Its New Minimum Age from the Dipole Turning Radius,''
		in \emph{YUCT}, Zenodo (2026).
		
		\bibitem{Whitmire1985}
		D. P. Whitmire, J. J. Matese, ``Periodic comet showers and the Nemesis hypothesis,''
		\emph{Nature} \textbf{313}, 36–38 (1985).
		
		\bibitem{Kirkpatrick2021}
		D. Kirkpatrick et al., ``The field brown dwarf population and the WISE survey,''
		\emph{Astrophysical Journal Supplement} \textbf{253}, 7 (2021).
		
		\bibitem{Korenaga2017}
		J. Korenaga, ``Earth’s thermal evolution, mantle convection, and Hadean onset of plate tectonics,''
		\emph{Earth and Planetary Science Letters} \textbf{145}, 334–348 (2017).
		
		\bibitem{Charnay2017}
		B. Charnay, G. Le Hir, F. Fluteau, F. Forget, D. C. Catling, ``A warm or a cold early Earth? New insights from a 3‑D climate‑carbon model,''
		\emph{Earth and Planetary Science Letters} \textbf{474}, 97–109 (2017).
		
		\bibitem{KrissansenTotton2019}
		J. Krissansen‑Totton, G. Arney, D. C. Catling, ``Ocean pH, atmospheric CO₂, and habitability of the early Earth,''
		\emph{AGU Fall Meeting} PP53B‑06 (2019).
		
		\bibitem{Arndt2012}
		N. T. Arndt, E. G. Nisbet, ``Processes on the Young Earth and the Habitats of Early Life,''
		\emph{Annual Review of Earth and Planetary Sciences} \textbf{40}, 521–549 (2012).
		
		\bibitem{Zircon2024}
		H. Gamaleldien et al., ``Four‑billion‑year‑old zircons reveal early fresh water and emerged continental crust,''
		\emph{Nature Geoscience} (2024). DOI: 10.1038/s41561‑024‑01450‑0
		
		\bibitem{PanspermiaReview}
		P. A. Bland, ``Panspermia: A viable hypothesis?''
		\emph{Astronomy \& Geophysics} \textbf{58}(5), 5.24–5.28 (2017).
		
		\bibitem{LunarMagmaOcean}
		M. Maurice, N. Tosi, S. Schwinger, D. Breuer, T. Kleine, ``The longevity of a lunar magma ocean: Implications for the Moon's thermal and magnetic evolution,''
		\emph{Journal of Geophysical Research: Planets} \textbf{125}, e2019JE006152 (2020). DOI: 10.1029/2019JE006152
		
		\bibitem{Chemolithoautotroph}
		Wikipedia, ``Chemolithoautotroph,''
		\url{https://en.wikipedia.org/wiki/Chemolithoautotroph} (accessed 2026).
		
		\bibitem{AbiogenesisReview}
		N. Kitadai, S. Maruyama, ``Onset of plate tectonics and the origin of life: A view from the Hadean,''
		\emph{Geoscience Frontiers} \textbf{9}(4), 1105–1122 (2018). DOI: 10.1016/j.gsf.2017.05.005
		
		\bibitem{Chemoautotroph}
		K. O. Stetter, ``Hyperthermophilic procaryotes,''
		\emph{FEMS Microbiology Reviews} \textbf{18}(2-3), 149–158 (1996).
		
		\bibitem{JanusEpimetheus}
		F. Namouni, H. Christophe, ``On the stability of co‑orbital motion of Janus and Epimetheus,''
		\emph{Astronomy \& Astrophysics} \textbf{434}, 1179–1186 (2005).
		
		\bibitem{Kearey2013}
		P. Kearey, K. A. Klepeis, F. J. Vine, \emph{Global Tectonics}, 3rd ed., Wiley-Blackwell (2013).
		
		\bibitem{Sahijpal2020}
		S. Sahijpal, V. Goyal, ``Thermal evolution of the early Moon,''
		\emph{Meteoritics \& Planetary Science} \textbf{53}(10), 2193–2210 (2018). DOI: 10.1111/maps.13119
		
		\bibitem{Colin2024}
		L. Colin, M. Maurice, N. Tosi, et al., ``Thermal evolution of the lunar magma ocean,''
		\emph{Earth and Planetary Science Letters} \textbf{648}, 119072 (2024). DOI: 10.1016/j.epsl.2024.119072
		
		\bibitem{Maurice2018}
		M. Maurice, N. Tosi, S. Schwinger, et al., ``A long-lived lunar magma ocean,''
		\emph{AGU Fall Meeting Abstracts}, P43C-06 (2018).
		
		\bibitem{Maurice2020}
		M. Maurice, N. Tosi, S. Schwinger, et al., ``A long-lived magma ocean on a young Moon,''
		\emph{Science Advances} \textbf{6}(28), eaba8949 (2020). DOI: 10.1126/sciadv.aba8949
		
		\bibitem{Zahnle2015}
		K. J. Zahnle, R. Lupu, D. C. Catling, ``The evolution of a steam atmosphere after a giant impact,''
		\emph{AGU Fall Meeting} P51A-1920 (2015).
		
		\bibitem{Zahnle2007}
		K. J. Zahnle, N. Arndt, C. Cockell, et al., ``Emergence of a habitable planet,''
		\emph{Space Science Reviews} \textbf{129}(1-3), 35–78 (2007). DOI: 10.1007/s11214-007-9225-z
		
		\bibitem{Feulner2012}
		G. Feulner, ``The faint young Sun problem,''
		\emph{Reviews of Geophysics} \textbf{50}(2), RG2006 (2012). DOI: 10.1029/2011RG000375
		
		\bibitem{Canup2014}
		R. M. Canup, ``Lunar-forming impacts: processes and alternatives,''
		\emph{Philosophical Transactions of the Royal Society A} \textbf{372}(2024), 20130175 (2014). DOI: 10.1098/rsta.2013.0175
		
	\end{thebibliography}
	
\end{document}